\documentclass[10pt,letterpaper]{article}
\usepackage[T1]{fontenc}
\usepackage{lmodern,amsmath,amssymb,amsthm,mathtools}
\usepackage[margin=0.82in]{geometry}
\usepackage{booktabs,longtable,array,microtype,xcolor,enumitem}
\usepackage{hyperref,xurl,fancyhdr}
\hypersetup{colorlinks=true,linkcolor=blue!45!black,urlcolor=blue!45!black,citecolor=blue!45!black,pdftitle={Reciprocal pooling across signed bases for Diophantine consistency},pdfauthor={AI-assisted research prepared for Vladimir Reshetnikov}}
\pagestyle{fancy}\fancyhf{}\fancyhead[L]{\small Reciprocal pooling and Diophantine consistency}\fancyhead[R]{\small ProveIt research contribution}\fancyfoot[C]{\thepage}
\setlength{\headheight}{13pt}
\setlength{\parskip}{4pt}\setlength{\parindent}{0pt}
\setlist{nosep,leftmargin=1.5em}
\newtheorem{theorem}{Theorem}[section]
\newtheorem{lemma}[theorem]{Lemma}
\newtheorem{corollary}[theorem]{Corollary}
\newtheorem{proposition}[theorem]{Proposition}
\theoremstyle{remark}\newtheorem{remark}[theorem]{Remark}
\newcommand{\Z}{\mathbb Z}\newcommand{\Q}{\mathbb Q}\newcommand{\N}{\mathbb N}
\newcommand{\Con}{\operatorname{Con}}
\newcommand{\sgn}{\operatorname{sgn}}
\newcommand{\Ht}{\operatorname{Ht}}
\newcommand{\src}{\textsf{S216}}
\title{Reciprocal pooling across signed bases\\for Diophantine consistency}
\author{AI-assisted research prepared for Vladimir Reshetnikov\\\small A contribution for the ProveIt repository}
\date{25 October 2026}

\begin{document}
\maketitle
\begin{abstract}
We give a uniform arithmetic compression theorem for finitely many
Pell-certified power requests with different, signed, or zero bases.
A single polynomial congruence witness replaces the separate witnesses
for all power blocks. The principal new point is a second integrality
argument: it rules out cancellation between positive and negative
reciprocal moduli without introducing orientation factors. Applied to an
explicit, calibrated numerical proof-calculus equation in the ProveIt
repository, it removes nine coordinates. Subsequent exact monic
eliminations give the fixed integer polynomials listed below.
For every delivered equation, ordinary mathematics formalizable in ZFC,
composed with the explicitly retained source calibration, proves that
absence of an integer zero is equivalent to syntactic $\Con(\mathrm{ZFC})$.
The complete finite polynomial definitions, source residuals, elimination
tables, executable circuits, and independent finite checkers are supplied.
We do not claim a globally simplest equation, a consistency proof, or a
new end-to-end proof-assistant certificate.
\end{abstract}

\section{Results and the meaning of simplicity}

Our principal recommendation is the fixed polynomial in \textbf{112 unrestricted integers of exact degree 24}. A lower-degree companion uses \textbf{119 integers at degree 16}; the smallest variable count in this release is \textbf{102 integers at degree 64}. Each has the required ZFC equivalence, with the dependency boundary described below.

\begin{center}

\begin{tabular}{rrrrr}

\toprule

Integer variables & Exact degree & Residuals & Height bound & Binary operations\\

\midrule

143 & 8 & 80 & $<2^{43}$ & 2590\\

119 & 16 & 56 & $<2^{79}$ & 2519\\

112 & 24 & 49 & $<2^{79}$ & 2501\\

107 & 32 & 44 & $<2^{113}$ & 2482\\

102 & 64 & 39 & $<2^{214}$ & 2437\\

\bottomrule

\end{tabular}

\end{center}

The 143-variable octic is a useful companion obtained with five singleton blocks and ordinary elimination; it does not need the new cross-base fusion theorem. The other four rows use that theorem. Operation counts describe the delivered literal evaluators and are not claimed optimal.

For comparison, package 152 gives 153 integers at degree 8, 125 at degree 16, 117 at degree 32, and 111 at degree 64 \cite{eliminate152}. The corresponding supplied coefficient-height bounds are $2^{47},2^{80},2^{144},2^{173}$. Our results lower the variable counts at those degrees; our degree-64 height bound $2^{214}$ is larger than the older bound $2^{173}$. Thus even this direct-family comparison has a coefficient tradeoff.

A fair local control isolates the new arithmetic gain. Apply exactly the principal form\textquotesingle s 95 core eliminations to the original source and also eliminate its three modulus coordinates and two already monic $K$ coordinates. This gives 116 integers, 53 residuals, and exact degree 24 without cross-base fusion. The new construction has 112 integers, 49 residuals, and the same degree. The four-coordinate saving in this matched comparison is attributable to the new fusion; the larger comparison against older packages also benefits from the improved source and generic elimination.


All variables in this report are independently quantified over $\Z$.
In particular, no nonnegativity condition on a free coordinate is silently
assumed. The coefficient data are fixed integers. The final polynomials
use only addition, subtraction and multiplication, with fixed finite
index sets. Exponentiation by a variable occurs in the \emph{proof of
correctness}, never in the definition of a delivered equation.

The main contribution is a reusable compression theorem, followed by an
explicit instantiation. It does not introduce a new formalization of set
theory. The selected source encodes the finite numerical calculus C222
directly, and its arithmetic calibration connects C222 refutations to
ZF and ZFC refutations. The present polynomial is therefore not a
simulation of K215 and does not instantiate Jones's universal polynomial.
Those distinctions concern the route used, not an assertion that their
other tradeoffs have become inferior.

There is no agreed scalar measure of ``simplest.'' Variable count, total
degree, coefficient height, number of monomials, straight-line-program
length, and the size of fixed coefficient data compete. The repository
already includes a quartic in 37 unrestricted integers, nine-positive-
variable equations of much higher degree, and ten-integer equations
using a different reduction \cite{corpus}. Our contribution concerns the
direct C222 route with small literal integer data. It does not replace
those constructions or establish a lower bound on any unrestricted
notion of simplicity.

\subsection{Comparison boundary and attribution}

The repository snapshot is
\begin{center}\small\texttt{ef0c534c1ae1e312c61a2a9a5a3f18f5516fd42b}.\end{center}
The corpus was inspected through package 172. The main README's summary
lags some delivered packages; our arithmetic source is the actual
pre-row-reduction 216-coordinate system in package 172, not that summary.
The exact source bytes have SHA-256
\begin{center}\scriptsize\texttt{e5e14bc0bbabb625f78ce57017fbdb336c4a65c5877f517142c399769992ff39}.\end{center}

Package 158 introduced pooling at one nonzero base. Package 168 combined
the two consecutive-base pools by a reciprocal sum, relying on their
common modulus sign, and discussed orientation weights as a possible
generalization \cite{pool158,pool168}. Package 172 supplied the present
arithmetic source, including its signed bootstrap, exact projections,
and ordinary C222 calibration \cite{source172}. Package 152 supplied the
general monic-elimination method used for the later degree/variable
tradeoffs \cite{eliminate152}. We credit all of these ingredients.

The new mathematical step is the unweighted pooling theorem below,
including mixed signs and zero bases, and its explicit composition with
the authenticated source. The search for elimination orders is a finite
optimization of that construction. No literature-wide priority claim is
made. Independent reviews here mean separately assigned AI-assisted
mathematical and program checks, not outside human refereeing.

\section{The arithmetic interface}

For an integer $A$, let the canonical signed Pell coordinates be defined
by
\[
 (A+\sqrt{A^2-1})^k=X_k(A)+Y_k(A)\sqrt{A^2-1},\qquad k\in\Z.
\]
This notation can equivalently be defined by integer recurrence, avoiding
any analytic choice of a square root. Put
\[
 M_A=\begin{pmatrix}A&A^2-1\\1&A\end{pmatrix},\qquad
 \begin{pmatrix}X_k(A)\\Y_k(A)\end{pmatrix}
 =M_A^k\begin{pmatrix}1\\0\end{pmatrix}.
\]
The determinant is one, so negative powers also have integer entries.
In particular $X_{-k}=X_k$ and $Y_{-k}=-Y_k$. Both signs of $A$ are
permitted throughout the proof.

Suppose that $s$ bases $n_1,\ldots,n_s$ support a total of $r$ requests.
For $i\in I_j$, the proposed request is $m_i=n_j^{k_i}$ and the retained
Pell pair is $(X_i,Y_i)=(X_{k_i}(A),Y_{k_i}(A))$. The sets $I_j$ are fixed
finite index sets. Define
\begin{align}
 T_j&=2An_j-n_j^2-1,\label{eq:T}\\
 \Delta_i&=m_i-X_i+AY_i,\notag\\
 H_j&=\sum_{i\in I_j}Y_i(m_i-X_i),\notag\\
 S_j&=\sum_{i\in I_j}(\Delta_i^2-Y_i^2),\notag\\
 E_j&=S_j-2n_jH_j.\label{eq:E}
\end{align}
In the literal source, $Z_i=AY_i$ is an already retained coordinate or
affine expression. Thus $\Delta_i=m_i-X_i+Z_i$ is affine, $H_j,S_j$ are
quadratic, and $E_j$ is at most cubic. The letters in these displays
are polynomial abbreviations unless explicitly declared as coordinates.

\begin{lemma}[Pell congruence for signed indices]\label{lem:pell}
If $n\ne0$ and $T=2An-n^2-1$, then $\gcd(n,T)=1$. Whenever $T\ne0$,
\[
 X_k(A)+(n-A)Y_k(A)\equiv n^k\pmod T\qquad(k\in\Z),
\]
where negative exponents on the right mean powers of the unit $n$
in $\Z/T\Z$.
\end{lemma}
\begin{proof}
The identity $T\equiv-1\pmod n$ proves coprimality. For the row vector
$\ell=(1,n-A)$, direct multiplication gives
\[
 \ell M_A=(n,An-1)\equiv n(1,n-A)=n\ell\pmod T,
\]
because $An-1-n(n-A)=T$. Both $M_A$ and $n$ are invertible modulo
$T$, so this relation iterates through every signed integer $k$.
Apply it to $(1,0)^{\mathsf T}$. No sign or parity assumption on $A$
or $k$ has been used.
\end{proof}

Write $\delta_i=m_i-X_i+(A-n_j)Y_i$ for $i\in I_j$. Expansion gives
the central identity
\begin{equation}
 \sum_{i\in I_j}\delta_i^2
 =E_j-T_j\sum_{i\in I_j}Y_i^2.\label{eq:normidentity}
\end{equation}
The symbolic checker verifies this identity using only the retained
$Z_i-AY_i$ equations, with all actual coordinate offsets included.

\section{Unweighted pooling across arbitrary signed bases}

\begin{theorem}[Two integrality steps]\label{thm:main}
Let $W\ge258$ and $|A|>W^W$, and put
$b=\lfloor\sqrt{W-1}\rfloor$. Suppose
\[
 n_j^2<W,\qquad m_i^2<W,\qquad k_i^2<W,
\]
and suppose the Pell pairs are canonical at their stated signed indices.
Assume also
\begin{equation}
 4r<W,\qquad
 (s+1)b+\left\lceil\frac{s+1}{2}\right\rceil+1\le W.
 \label{eq:budget}
\end{equation}
The following conditions are equivalent:
\begin{enumerate}
\item $T_j\mid E_j$ for every $j$;
\item every request at a nonzero base is the exact rational equality
      $m_i=n_j^{k_i}$;
\item there is $z\in\Z$ such that
\begin{equation}
 \boxed{\quad
 R=\sum_{j=1}^{s}E_j\prod_{h\ne j}T_h
       -z\prod_{h=1}^{s}T_h=0.\quad}\label{eq:pooled}
\end{equation}
\end{enumerate}
For $|n_j|\ge2$, the same conditions force $k_i\ge0$.
Zero bases, unit bases, repeated bases, and noncoprime moduli are allowed.
In particular, the two budget inequalities hold uniformly for $s\le14$, $r\le64$.
\end{theorem}

\begin{proof}
First note that every $T_j$ is nonzero. If $n_j=0$, then $T_j=-1$.
Otherwise, since $1\le|n_j|\le b$,
\begin{equation}
 |T_j|\ge |n_j|(2|A|-W)>|n_j|W^W,
 \qquad \sgn(n_jT_j)=\sgn(A).\label{eq:denombound}
\end{equation}
Indeed $|n_j|+1/|n_j|\le b+1\le W$, and $|A|>W^W>W$.

Assume \eqref{eq:pooled}. Divide by the nonzero product of moduli and
use \eqref{eq:normidentity}. Then
$\sum_{j,i\in I_j}\delta_i^2/T_j$ is an integer. Terms at zero bases
are already integers, because their denominators equal $-1$. Remove
them from this rational sum. If no bases remain, every required
divisibility is automatic. Otherwise let $J=\{j:n_j\ne0\}$ and set
\[
 L=\prod_{j\in J}n_j,\qquad C=L^b\ne0,
 \qquad
 d_i=Cm_i-\frac{C}{n_j^b}n_j^{b+k_i}\quad(i\in I_j,\ j\in J).
\]
These are integers: $b+k_i\ge0$ and $C/n_j^b$ is a product of
integer powers. Lemma~\ref{lem:pell} implies
$C\delta_i\equiv d_i\pmod{T_j}$, including negative indices.
Consequently, with
\[
 D_j^*=\sum_{i\in I_j}d_i^2,\qquad D^*=\sum_{j\in J}D_j^*,
\]
the rational number $R_0=\sum_{j\in J}D_j^*/T_j$ differs by an integer
from $C^2\sum\delta_i^2/T_j$. Hence $R_0\in\Z$.

Every input has absolute value at most $b$. If $s'=|J|\le s$, then
\[
 |d_i|\le b^{s'b+1}+b^{(s'+1)b}\le2b^{(s+1)b},\qquad
 D^*\le4r b^{2(s+1)b}.
\]
The stronger estimate needed momentarily is
\begin{align}
 2|L|D^*
 &\le8r b^{2(s+1)b+s}
 \le4r b^{2(s+1)b+s+1}\notag\\
 &<W^{(s+1)b+\lceil(s+1)/2\rceil+1}
 \le W^W.\label{eq:strongbudget}
\end{align}
Here $2\le b$, $b^2<W$, $4r<W$, and \eqref{eq:budget} are sufficient.
In particular $D^*<W^W$. Thus \eqref{eq:denombound} gives
$|R_0|\le D^*/(2|A|-W)<1$. Since $R_0$ is an integer, $R_0=0$.

This is not yet enough: denominators can have opposite signs. The second
integrality step uses the exact identity
\begin{equation}
 \frac{2A}{T_j}=\frac1{n_j}
       +\frac{n_j^2+1}{n_jT_j}.\label{eq:split}
\end{equation}
Multiplying $R_0=0$ by $2A$ gives $U+V=0$, where
\[
 U=\sum_{j\in J}\frac{D_j^*}{n_j},\qquad
 V=\sum_{j\in J}\frac{D_j^*(n_j^2+1)}{n_jT_j}.
\]
Now $LU\in\Z$, so $LV=-LU\in\Z$. On the other hand,
\[
 \frac{n_j^2+1}{|n_jT_j|}
 \le\frac{1+n_j^{-2}}{2|A|-W}
 \le\frac2{2|A|-W},
\]
and \eqref{eq:strongbudget} yields
$|LV|\le2|L|D^*/(2|A|-W)<1$. Therefore $LV=0$ and $V=0$.
Every nonzero summand of $V$ has the same sign, namely $\sgn(A)$.
All $D_j^*$ and all $d_i$ must vanish. Dividing $d_i=0$ in $\Q$ by
$C$ proves $m_i=n_j^{k_i}$. For a nonunit base and a negative index,
this would imply $m_i n_j^{-k_i}=1$, which is impossible in $\Z$.

Conversely, exact powers at nonzero bases and Lemma~\ref{lem:pell}
give $T_j\mid\delta_i$. Identity~\eqref{eq:normidentity} then gives
$T_j\mid E_j$. At a zero base this divisibility is automatic. Finally,
if every $T_j\mid E_j$, choose $z=\sum_j E_j/T_j\in\Z$ to obtain
\eqref{eq:pooled}.

For the uniform numerical corollary, $b\ge16$ and
\[
 (s+1)b+\lceil(s+1)/2\rceil+1\le15b+9\le b^2+1\le W
\]
when $s\le14$. Also $4r\le256<W$ when $r\le64$.
\end{proof}

\subsection{Why the second step matters}

A single small-integer argument only proves $R_0=0$. It does not prevent
cancellation between moduli of opposite signs. The decomposition
\eqref{eq:split} separates a rational leading term with controlled
denominator from a smaller correction whose sign is common. This is the
additional mechanism that permits removal of the orientation factors.

The height assumption is essential. Set $A=10$ and let every index be
zero, so every canonical pair is $(1,0)$. At base $+1$ take proposed
output $4$, giving $E_+=9,T_+=18$. At base $-1$ take outputs $4,2,2$,
giving $E_-=11,T_-=-22$. Then
\[
 E_+/T_++E_-/T_- =\tfrac12-\tfrac12=0,
\]
although neither divisibility holds. This example violates the large-
center hypothesis. It is included in the independent audit as a concrete
reason not to delete that hypothesis.

\subsection{Integral witness recovery}

The retained source uses two kinds of power block. An individual request
has private coordinates $T,Q$ and rows
\[
 T-(2An-n^2-1)=0,\qquad \delta+TQ=0.
\]
After applying the theorem, recover
\[
 T=2An-n^2-1,\qquad Q=-\delta/T\in\Z.
\]
At a nonzero base, exact powers guarantee the required divisibility.
At a zero base, $T=-1$ guarantees it without interpreting $0^k$.

A same-base block has private $K,L$ and rows
\begin{equation}
 K-2H-(2A-n)L=0,\qquad S-nK+L=0.\label{eq:oldpool}
\end{equation}
Restore
\[
 L=E/T,\qquad K=2H+(2A-n)L.
\]
Both are integers and satisfy \eqref{eq:oldpool}. For $n=0$, simply
$L=-E$. Thus the new equation restores the full old polynomial system
before the old semantic bootstrap is invoked. No circular argument
about nonzero bases is introduced.

These maps preserve the set of retained tuples that admit witnesses.
They are not polynomial coordinate bijections on all of affine space:
division is used only in the proof of integral recovery on the zero set.
The emitted polynomial itself contains no division.

\section{The explicit calibrated source and its transformation}

\subsection{The five blocks}

The literal source $\src$ has 216 unrestricted integer coordinates
$x_0,\ldots,x_{215}$ and 153 quadratic residuals
$r_0,\ldots,r_{152}$. Appendix~\ref{app:rows} prints every residual.
The arithmetic roles used here are
\[
 A=1+x_{133},\quad t=x_0,\quad F=x_3,\quad B=x_4,
 \quad BH=x_{185},\quad q=1+4x_{209}.
\]
The five bases, in their fixed order, are
\begin{equation}
 (n_0,n_1,n_2,n_3,n_4)=(t,B,BH,q+1,q).\label{eq:fivebases}
\end{equation}
The first three blocks each have one request. The fourth has eleven and
the fifth has three, for a total of seventeen. The complete request table
appears in Appendix~\ref{app:requests}. Repeated use of a Pell pair at
the same index does not create a new coordinate.

Delete exactly the ten rows
\begin{equation}
 \mathcal D=\{97,98,99,100,101,102,149,150,151,152\}\label{eq:deletedrows}
\end{equation}
and the ten coordinates
\begin{equation}
 \mathcal V=\{140,141,142,143,144,145,212,213,214,215\}.
 \label{eq:deletedvars}
\end{equation}
The first six coordinates are the three individual moduli and quotients;
the last four are the two $K,L$ pairs. They occur nowhere in the retained
143 rows. Their exact private occurrences and their signs are checked.

Define the five $T_j,H_j,S_j,E_j$ from the complete request table and
\eqref{eq:T}--\eqref{eq:E}, and let $R$ be \eqref{eq:pooled} with
$s=5$. The first explicit new polynomial is
\begin{equation}
 \boxed{\quad P_{207}=\sum_{i\notin\mathcal D}r_i^2+R^2.\quad}
 \label{eq:P207}
\end{equation}
Its free coordinates are the 206 old coordinates outside $\mathcal V$
and the new integer $z$. For software, the retained old indices are
renumbered in increasing order, followed by $z$; the exact map is stored.

Because every base and $A$ is affine, each $T_j$ is quadratic and each
$E_j$ is at most cubic. Thus $R$ has degree at most 11. Its coefficient
of the private variable $z$ is $-\prod_jT_j$, a polynomial of degree
exactly ten. Each $T_j$ has a nonzero quadratic homogeneous part, and
$\Z[\mathbf x]$ is an integral domain. Therefore $\deg R=11$ and
$\deg P_{207}=22$ exactly. Taking sums of squares introduces no
cancellation of the highest nonzero homogeneous square.

\subsection{Why the retained hypotheses still hold}

The source proof establishes canonical signed Pell pairs and the large
center before interpreting the power congruences or recovering front
signs. The actual frame is
\begin{equation}
 W=2+\sum_{a\in\mathcal A}a^2,\qquad |\mathcal A|=34,
 \label{eq:frame}
\end{equation}
where the complete affine-atom list is supplied in
\texttt{inputs/interface216.json}. The independent checker expands
\eqref{eq:frame} and matches source row 129 exactly. Every actual base,
output and index of all seventeen requests is one of these atoms.
The $t$ atom and retained horizon equation give $t\ge16$, hence
$W\ge258$. The retained scale and semantic-anchor argument gives
$|A|>W^W$ \cite{source172}.

All fourteen $Z_i=AY_i$ equations, fourteen Pell norm equations, the
complete signed-orientation equation, the scale, and the frame remain.
The checker also reconstructs source row 142 as the orientation identity
minus the fourteen retained norm identities. Three delicate conventions
are checked literally: $A$ is stored with offset $+1$; pair 3's ordinate
has offset $-1$; pair 12's $Z$ has offset $+1$. A variable name alone
is never taken as evidence of its arithmetic value.

Here $s=5,r=17$, so $4r=68<258$ and
$6b+4\le b^2+1\le W$. Theorem~\ref{thm:main} applies to every zero
of the new retained system. It restores all ten deleted witnesses and
all ten source rows. Conversely, every source zero supplies an integer
$z$. Thus $P_{207}$ and $\src$ have equivalent integer-zero existence.

\subsection{The logical endpoint}

The ordinary source calibration proves
\begin{align*}
 \exists\mathbf x\in\Z^{216}\ (r_i(\mathbf x)=0\text{ for all }i)
 &\quad\Longleftrightarrow\quad
 \text{there is a finite C222 refutation}\\
 &\quad\Longleftrightarrow\quad
 \text{there is a finite ZF refutation}\\
 &\quad\Longleftrightarrow\quad
 \text{there is a finite ZFC refutation}.
\end{align*}
The named logical interfaces are
\begin{quote}\small\ttfamily
KEndpoint.c222\_iff\_zf\_refutation\\
ConstructibleChoice.zf\_refutation\_iff\_zfc\_refutation.
\end{quote}
Their provenance and source files are retained with the parent arithmetic
proofs. The appearance of the name \texttt{KEndpoint} in an inherited
interface is not a Turing-machine simulation in the present polynomial.

All new steps are finite integer algebra, elementary bounds, recurrence
arguments, and explicit finite-witness transformations. They are
formalizable in ZFC and do not assume $\Con(\mathrm{ZFC})$.
Composing with the identified ordinary source calibration gives
\begin{equation}
 \mathrm{ZFC}\vdash
 \left(\forall\mathbf u\in\Z^{207},\ P_{207}(\mathbf u)\ne0\right)
 \longleftrightarrow\Con(\mathrm{ZFC}).\label{eq:logic}
\end{equation}
Here $\Con(\mathrm{ZFC})$ refers to absence of a finite refutation in
the effective presentation fixed by the inherited calibration. No
semantic soundness assumption on ZFC is needed to state or prove the
equivalence. The report does not prove either side of it.

\section{Removing defined coordinates}

\begin{lemma}[Acyclic monic elimination]\label{lem:eliminate}
Suppose selected residuals have the form
$r_j=\varepsilon x_v+f$, where $\varepsilon\in\{1,-1\}$ and $f$
does not contain $x_v$. If their dependency graph is acyclic, replace
every selected $x_v$ by the recursively defined polynomial
$-\varepsilon f$, omit the selected residuals, and retain every other
residual after substitution. This gives a polynomial bijection between
the zero sets in the remaining and original coordinates.
\end{lemma}
\begin{proof}
Resolve the definitions in a topological order. Each gives one unique
integer coordinate for any values of its already resolved arguments.
The omitted row then vanishes identically. The retained rows become
exactly the substituted rows. Forward restriction and reverse polynomial
restoration are mutual inverses on the zero sets.
\end{proof}

This lemma is a standard exact elimination step and was already applied
extensively in package 152. The contribution here is its composition with
the new power pooling and the explicit selected pivot tables. Every
pivot is checked against a retained \emph{original} source row: its
coefficient is exactly $+1$ or $-1$, it has no other occurrence in that
row, and no private power coordinate is referenced. The compiler checks
acyclicity and recursively substitutes the definitions. No rational
Gaussian-elimination step is used to remove an integer coordinate.

The principal smaller polynomials are fully defined by
\eqref{eq:P207}, or by its specified block partition, followed by the
finite pivot lists in Appendix~\ref{app:pivots}. A pivot $(i,v)$ always
means: solve the printed original row $r_i$ for $x_v$ with its unit
coefficient. Resolve all such definitions recursively, remove their
rows, and sum the squares of the remaining substituted residuals and
the specified pooled residuals. This is a complete polynomial recipe
with fixed integer coefficients; it does not call an unspecified
Diophantine compiler.

\subsection{Partitioning blocks to control degree}

The five blocks may be partitioned into fixed subsets. Apply
Theorem~\ref{thm:main} separately to each subset, introducing one $z$
per part. The unchanged square-frame budget is sufficient for every
part. Before monic eliminations, the selected partitions have these
exact counts:
\begin{center}
\begin{tabular}{lrrrr}
\toprule
Block sizes & Integer variables & Residuals & Degree & Height bound\\
\midrule
$5$ &207&144&22&$<2^{46}$\\
$3+2$ &208&145&14&$<2^{29}$\\
$2+2+1$ &209&146&10&$<2^{25}$\\
$1+1+1+1+1$ &211&148&6&$<2^{24}$\\
\bottomrule
\end{tabular}
\end{center}
Even a singleton part can replace its two old private coordinates by
one witness, at a degree cost. These rows are intermediate constructions,
not a claim that their raw variable counts improve the repository's
already eliminated high-degree equations.

\subsection{Exact degree, liveness, and coefficient size}

The emitted straight-line programs are complete, topologically ordered
arrays with integer leaves and binary $+,-,\times$ gates. Their inputs
are exactly the declared free integer coordinates. Sharing a node is
only an evaluator optimization and adds no mathematical variable.

Degrees are certified in two steps. Propagating syntactic degrees gives
an upper bound. Propagating homogeneous leading values at a supplied
integer point gives a nonzero value in that upper degree, proving
equality. The independent checkers recompute these witnesses. Nonzero
partial derivatives at explicit finite-field points certify that every
declared coordinate really occurs in the final polynomial; syntactic
appearance alone would not prove this after cancellation.

We have not expanded the largest output polynomials, so exact monomial
counts and exact expanded coefficient heights are not claimed. A useful
rigorous bound is nevertheless available. Assign norm bound $1$ to each
variable, $|c|$ to an integer constant, sum bounds at addition or
subtraction, and multiply bounds at multiplication. These rules bound
the coefficient $\ell^1$ norm, hence every expanded coefficient. If the
resulting positive integer has bit length $B$, then $\Ht(P)<2^B$.
The reported bounds use the complete emitted circuits, including the
sum of residual squares. The largest literal integer in those circuits
is 60; this must not be confused with their expanded coefficient height.

\section{Verification, reproducibility, and remaining limits}

The release contains two independent mathematical reviews of the
mixed-sign pooling argument. The first supplies the uniform theorem and
the signed Pell congruence proof; the second separately checks the
five-base/seventeen-request bounds, zero and unit bases, and the
small-center counterexample. The independent finite checks cover:
\begin{itemize}
\item the immutable source hash, all seventeen request ports, the exact
      34-atom frame, all fourteen Pell pairs, offsets, and every private
      occurrence of the ten removed coordinates;
\item the new $T,H,S,E$ expressions, the error-square identity, pool
      partitions, the complete retained residual list, and emitted circuits;
\item every monic pivot, its integer inverse definition, acyclicity,
      the full substituted circuit, exact degree and variable liveness;
\item arithmetic samples with both signs of $A$, mixed and zero bases,
      units with negative indices, false outputs, and impossible negative
      nonunit powers, together with deliberate corruptions;
\item deterministic reconstruction in a fresh directory and both ordinary
      and optimized Python execution, so disabled \texttt{assert}
      statements cannot silently remove the validation.
\end{itemize}

The restored package-172 source tree and its finite arithmetic ancestry
were also replayed. Its recorded receipts distinguish that replay from
a new Lean/Rocq build, which was not performed. The full ordinary proof
of the inherited calibration is a dependency; finite checks do not turn
it into a proof-assistant certificate. Similarly, numerical tests of
the new arithmetic theorem support implementation review but do not
replace the quantified proof in Section~3.

Run the release verifier from any working directory using an absolute
path to \texttt{verify.py}. The checker requires Python 3.10 or later
and its standard library. \texttt{evaluate.py} evaluates a selected
literal circuit on one supplied tuple, exactly or modulo a chosen
integer. It performs no root search. A modular zero has no asserted
relationship to a ZFC refutation.

The verification record in \texttt{verification/} identifies exactly
what was run. \texttt{SOURCES.json} and \texttt{NOTICE.md} distinguish
new results from inherited source files. The archive's manifest covers
the final release bytes. The \LaTeX{} sources and all generated inputs
are present; an ordinary \texttt{pdflatex} installation rebuilds this
report. Rebuilding may change PDF metadata without changing its content.

\section{Further research questions}

\paragraph{Can the two-scale argument reduce more shared certificates?}
The proof exploits the particular expansion
$2A/T=1/n+(n^2+1)/(nT)$. More generally, one can seek denominator
families with a rational leading term of bounded denominator and a
correction of one sign. A useful extension would state precise height
conditions and produce a polynomial with a smaller aggregate degree.
Merely taking an arbitrary reciprocal sum is invalid, as the explicit
small-center counterexample demonstrates.

\paragraph{How far can partition-aware elimination go?}
The selected orders are finite successful searches, not optima. The
choice of block partition interacts with which bases, moduli, and Pell
products may be inlined while keeping a target degree. A joint search
could use exact degree certificates during optimization rather than
only safe weighted-degree upper bounds. Any new point should be compared
against the same source with the fusion disabled, so the contribution
of the arithmetic theorem remains measurable.

\paragraph{Can the quotient witness be carried by another coordinate?}
The new aggregate adds one unrestricted integer per part. If a retained
coordinate has a demonstrably free cofinal choice, a polynomial carrier
might encode this quotient without adding a new coordinate. Such a
proposal needs both an integral soundness recovery and a constructive
completeness argument. No such carrier is claimed here.

\paragraph{Can a comparable gain be obtained at degree four?}
The present compression deliberately spends degree. Returning to a
quartic requires lifting products to new coordinates, which may erase
the variable savings. A useful objective is a structural identity that
retains the cancellation argument but realizes its conclusion with
fewer quadratic constraints than the original five blocks. The current
result supplies no lower bound ruling this out.

\paragraph{Can the calibration and the general theorem be formalized together?}
A focused formal development could first prove the rational two-
integrality lemma, then its signed-Pell specialization, and finally the
literal source transformation. That is smaller in scope than rebuilding
an entire universal Diophantine theorem. The source arithmetic compiler
and object-language ZFC formalization boundary would still need to be
tracked explicitly. We have not represented ordinary formalizability
as a completed kernel check.

\paragraph{What is the best presentation under a declared simplicity budget?}
An exact comparison should fix the domain and charge for coefficient
data, degree, and either expanded support or a specified arithmetic
circuit. The direct route here avoids a large native coefficient tuple
but uses more variables than the repository's other routes. Proving a
global minimum is far beyond the current evidence. A reproducible
Pareto catalog with matched controls would be a concrete next result.

\appendix
\section{Complete request interface}\label{app:requests}
The following table is literal data, not an informal description.
The source coordinate numbers are used throughout. The quantities
$X=D+Z$, $n$, $m$ and $k$ supply the definitions in Section~2.
The final request uses the affine expression
\[
 J=5x_8-3x_9+4x_{19}+2x_{22}+2x_{26}+2x_{28}+2x_{30}
 +2x_{32}+4x_{34}+2x_{35}+2x_{38}+2x_{41}+4x_{60}+2x_{61}.
\]
\begin{center}\small
\begin{tabular}{rrrrlll}
\toprule
Request & Block & $n$ & $m$ & $k$ & $D=X-AY$ & $(Y,Z)$\\
\midrule
0 & 0 & $x_{0}$ & $x_{3}$ & $x_{0}$ & $x_{137}$ & $(x_{138},x_{139})$\\
1 & 1 & $x_{4}$ & $x_{17}$ & $x_{1}$ & $x_{173}$ & $(x_{174},1+x_{175})$\\
2 & 2 & $x_{185}$ & $1+4x_{209}$ & $x_{1}$ & $x_{173}$ & $(x_{174},1+x_{175})$\\
3 & 3 & $2+4x_{209}$ & $x_{94}$ & $x_{0}+x_{3}$ & $x_{146}$ & $(x_{147},x_{148})$\\
4 & 4 & $1+4x_{209}$ & $x_{95}$ & $x_{0}$ & $x_{137}$ & $(x_{138},x_{139})$\\
5 & 3 & $2+4x_{209}$ & $x_{99}$ & $-1+x_{8}$ & $x_{149}$ & $(-1+x_{150},x_{151})$\\
7 & 3 & $2+4x_{209}$ & $x_{102}$ & $x_{9}$ & $x_{152}$ & $(x_{153},x_{154})$\\
8 & 3 & $2+4x_{209}$ & $x_{105}$ & $x_{10}$ & $x_{155}$ & $(x_{156},x_{157})$\\
9 & 3 & $2+4x_{209}$ & $x_{108}$ & $x_{11}$ & $x_{158}$ & $(x_{159},x_{160})$\\
10 & 3 & $2+4x_{209}$ & $x_{111}$ & $x_{12}$ & $x_{161}$ & $(x_{162},x_{163})$\\
11 & 4 & $1+4x_{209}$ & $x_{112}$ & $x_{2}$ & $x_{164}$ & $(x_{165},x_{166})$\\
12 & 3 & $2+4x_{209}$ & $x_{115}$ & $x_{13}$ & $x_{167}$ & $(x_{168},x_{169})$\\
13 & 3 & $2+4x_{209}$ & $x_{118}$ & $x_{14}$ & $x_{170}$ & $(x_{171},x_{172})$\\
14 & 4 & $1+4x_{209}$ & $x_{119}$ & $x_{1}$ & $x_{173}$ & $(x_{174},1+x_{175})$\\
15 & 3 & $2+4x_{209}$ & $x_{123}$ & $x_{15}$ & $x_{176}$ & $(x_{177},x_{178})$\\
16 & 3 & $2+4x_{209}$ & $x_{126}$ & $x_{16}$ & $x_{179}$ & $(x_{180},x_{181})$\\
17 & 3 & $2+4x_{209}$ & $x_{129}$ & $J$ & $x_{182}$ & $(x_{183},x_{184})$\\
\bottomrule
\end{tabular}
\end{center}


\section{Every quadratic source residual}\label{app:rows}
Set each of the following 153 polynomials equal to zero to recover the
literal source $\src$. Row numbers are zero-based. The rows in
\eqref{eq:deletedrows} are removed only by the proved pooling
transformation. No other row is silently dropped. These formulas,
the request table and the pivot tables specify every final polynomial
without requiring external numerical data.
\begingroup\small\allowdisplaybreaks[4]
% BEGIN_RESIDUAL 0
\begin{align*}
r_{0}={}&x_{0}-x_{0}x_{2}-x_{2}+x_{7}
\end{align*}
% END_RESIDUAL 0

% BEGIN_RESIDUAL 1
\begin{align*}
r_{1}={}&-1+x_{1}-x_{2}x_{5}-x_{3}
\end{align*}
% END_RESIDUAL 1

% BEGIN_RESIDUAL 2
\begin{align*}
r_{2}={}&-1+x_{1}-x_{6}x_{7}
\end{align*}
% END_RESIDUAL 2

% BEGIN_RESIDUAL 3
\begin{align*}
r_{3}={}&1+x_{3}-x_{4}x_{18}+x_{16}
\end{align*}
% END_RESIDUAL 3

% BEGIN_RESIDUAL 4
\begin{align*}
r_{4}={}&-1+x_{8}-x_{8}^{2}-2x_{8}x_{9}-x_{9}-x_{9}^{2}+x_{19}
\end{align*}
% END_RESIDUAL 4

% BEGIN_RESIDUAL 5
\begin{align*}
r_{5}={}&-5-20x_{8}-16x_{8}^{2}-16x_{8}x_{9}-16x_{8}x_{19}-10x_{9}-4x_{9}^{2}-8x_{9}x_{19}-10x_{19}-4x_{19}^{2}+x_{20}
\end{align*}
% END_RESIDUAL 5

% BEGIN_RESIDUAL 6
\begin{align*}
r_{6}={}&-2-x_{8}-x_{8}^{2}-4x_{8}x_{20}-4x_{20}-4x_{20}^{2}+x_{21}
\end{align*}
% END_RESIDUAL 6

% BEGIN_RESIDUAL 7
\begin{align*}
r_{7}={}&1+2x_{21}-4x_{21}^{2}+x_{22}
\end{align*}
% END_RESIDUAL 7

% BEGIN_RESIDUAL 8
\begin{align*}
r_{8}={}&-2-x_{9}^{2}-2x_{9}x_{10}-x_{10}-x_{10}^{2}+x_{23}
\end{align*}
% END_RESIDUAL 8

% BEGIN_RESIDUAL 9
\begin{align*}
r_{9}={}&-1+x_{8}-x_{8}^{2}-2x_{8}x_{10}-x_{9}-x_{10}^{2}+x_{24}
\end{align*}
% END_RESIDUAL 9

% BEGIN_RESIDUAL 10
\begin{align*}
r_{10}={}&8-11x_{8}-49x_{8}^{2}+10x_{8}x_{9}-44x_{8}x_{10}-32x_{8}x_{19}+4x_{8}x_{20}-4x_{8}x_{23}-32x_{8}x_{24}+25x_{9}+3x_{9}^{2}\\
&+8x_{9}x_{10}-8x_{9}x_{23}-15x_{10}-12x_{10}^{2}-16x_{10}x_{19}+8x_{10}x_{23}-16x_{10}x_{24}-9x_{19}-4x_{19}^{2}-8x_{19}x_{24}\\
&+4x_{20}+4x_{20}^{2}-x_{21}+12x_{23}-4x_{23}^{2}-8x_{24}-4x_{24}^{2}+x_{25}
\end{align*}
% END_RESIDUAL 10

% BEGIN_RESIDUAL 11
\begin{align*}
r_{11}={}&-18+14x_{8}-2x_{8}^{2}-8x_{8}x_{9}+8x_{8}x_{10}-8x_{8}x_{23}+24x_{9}-8x_{9}^{2}+16x_{9}x_{10}-16x_{9}x_{23}-30x_{10}-44x_{10}^{2}\\
&+16x_{10}x_{23}+24x_{10}x_{25}+24x_{23}-8x_{23}^{2}-4x_{25}^{2}+x_{26}
\end{align*}
% END_RESIDUAL 11

% BEGIN_RESIDUAL 12
\begin{align*}
r_{12}={}&-6-8x_{8}-4x_{8}^{2}-8x_{8}x_{9}-8x_{8}x_{10}+8x_{8}x_{23}+8x_{8}x_{24}-6x_{9}-4x_{9}^{2}-8x_{9}x_{10}+8x_{9}x_{23}+8x_{9}x_{24}\\
&-5x_{10}-4x_{10}^{2}+8x_{10}x_{23}+8x_{10}x_{24}-x_{19}+8x_{23}-4x_{23}^{2}-8x_{23}x_{24}+6x_{24}-4x_{24}^{2}+x_{27}
\end{align*}
% END_RESIDUAL 12

% BEGIN_RESIDUAL 13
\begin{align*}
r_{13}={}&1-4x_{8}+2x_{9}-4x_{9}^{2}-8x_{9}x_{10}-8x_{9}x_{27}+4x_{10}-4x_{10}^{2}-8x_{10}x_{27}-2x_{19}+4x_{27}-4x_{27}^{2}+x_{28}
\end{align*}
% END_RESIDUAL 13

% BEGIN_RESIDUAL 14
\begin{align*}
r_{14}={}&35-16x_{8}+2x_{8}^{2}+18x_{8}x_{9}-12x_{8}x_{10}-20x_{8}x_{23}-32x_{8}x_{24}-8x_{9}-14x_{9}^{2}-24x_{9}x_{10}-40x_{9}x_{23}\\
&-12x_{9}x_{24}-26x_{10}-24x_{10}^{2}-40x_{10}x_{23}-32x_{10}x_{24}-x_{19}-40x_{23}+4x_{23}^{2}+16x_{23}x_{24}-28x_{24}+4x_{24}^{2}\\
&-2x_{27}+x_{29}
\end{align*}
% END_RESIDUAL 14

% BEGIN_RESIDUAL 15
\begin{align*}
r_{15}={}&8+34x_{8}+18x_{8}^{2}-8x_{8}x_{9}+24x_{8}x_{10}+8x_{8}x_{23}+32x_{8}x_{24}-50x_{9}+10x_{9}^{2}-32x_{9}x_{10}+16x_{9}x_{23}\\
&-8x_{9}x_{24}+60x_{10}-4x_{10}^{2}-16x_{10}x_{23}+32x_{10}x_{24}-24x_{10}x_{25}-32x_{10}x_{29}-24x_{23}+8x_{23}^{2}+32x_{24}\\
&+8x_{24}^{2}+4x_{25}^{2}-x_{26}-2x_{29}-4x_{29}^{2}+x_{30}
\end{align*}
% END_RESIDUAL 15

% BEGIN_RESIDUAL 16
\begin{align*}
r_{16}={}&6-2x_{8}+x_{8}^{2}+4x_{8}x_{20}-47x_{9}-28x_{9}^{2}+16x_{9}x_{10}+24x_{9}x_{19}+24x_{9}x_{23}+7x_{10}+8x_{10}^{2}+16x_{19}\\
&-4x_{19}^{2}-8x_{19}x_{23}+4x_{20}+4x_{20}^{2}-x_{21}+6x_{23}-4x_{23}^{2}-x_{24}+x_{31}
\end{align*}
% END_RESIDUAL 16

% BEGIN_RESIDUAL 17
\begin{align*}
r_{17}={}&-3-6x_{8}+2x_{9}-2x_{10}-2x_{24}-4x_{31}-4x_{31}^{2}+x_{32}
\end{align*}
% END_RESIDUAL 17

% BEGIN_RESIDUAL 18
\begin{align*}
r_{18}={}&16+30x_{8}+14x_{8}^{2}-2x_{8}x_{9}-24x_{8}x_{10}-16x_{8}x_{19}+60x_{8}x_{20}-16x_{8}x_{27}+19x_{9}-6x_{9}x_{10}-24x_{9}x_{19}\\
&+32x_{9}x_{20}-12x_{10}+6x_{10}^{2}-12x_{10}x_{19}+12x_{10}x_{27}+3x_{19}-8x_{19}^{2}+32x_{19}x_{20}-8x_{19}x_{27}+15x_{20}\\
&-4x_{20}^{2}-x_{23}-2x_{24}+4x_{27}^{2}+x_{33}+8x_{36}-8x_{36}^{2}+2x_{37}-x_{42}
\end{align*}
% END_RESIDUAL 18

% BEGIN_RESIDUAL 19
\begin{align*}
r_{19}={}&1+2x_{33}-4x_{33}^{2}+x_{34}
\end{align*}
% END_RESIDUAL 19

% BEGIN_RESIDUAL 20
\begin{align*}
r_{20}={}&-x_{9}^{2}-4x_{9}x_{34}-4x_{34}^{2}+x_{35}
\end{align*}
% END_RESIDUAL 20

% BEGIN_RESIDUAL 21
\begin{align*}
r_{21}={}&-1-15x_{8}-17x_{8}^{2}+30x_{8}x_{9}-32x_{8}x_{19}+7x_{9}-13x_{9}^{2}+8x_{9}x_{10}+32x_{9}x_{19}-x_{10}-14x_{19}-16x_{19}^{2}\\
&-4x_{23}+x_{24}+x_{36}+2x_{71}
\end{align*}
% END_RESIDUAL 21

% BEGIN_RESIDUAL 22
\begin{align*}
r_{22}={}&-4-3x_{8}-7x_{9}-4x_{10}-2x_{23}-x_{24}+4x_{36}-4x_{36}^{2}+x_{37}
\end{align*}
% END_RESIDUAL 22

% BEGIN_RESIDUAL 23
\begin{align*}
r_{23}={}&-6x_{8}+2x_{9}+2x_{10}-4x_{10}^{2}-8x_{10}x_{37}-2x_{24}+4x_{37}-4x_{37}^{2}+x_{38}
\end{align*}
% END_RESIDUAL 23

% BEGIN_RESIDUAL 24
\begin{align*}
r_{24}={}&-x_{10}-x_{10}^{2}-2x_{10}x_{11}-x_{11}^{2}+x_{39}
\end{align*}
% END_RESIDUAL 24

% BEGIN_RESIDUAL 25
\begin{align*}
r_{25}={}&-31+2x_{8}-x_{8}^{2}-4x_{8}x_{39}-22x_{9}+16x_{9}^{2}-32x_{9}x_{10}+32x_{9}x_{23}-42x_{10}-20x_{10}^{2}-8x_{10}x_{23}\\
&+24x_{10}x_{24}+12x_{23}^{2}-8x_{23}x_{24}+22x_{24}-4x_{24}^{2}+2x_{39}-4x_{39}^{2}+x_{40}-4x_{43}
\end{align*}
% END_RESIDUAL 25

% BEGIN_RESIDUAL 26
\begin{align*}
r_{26}={}&4x_{8}-2x_{8}^{2}-8x_{8}x_{39}-16x_{10}^{2}-16x_{10}x_{40}+4x_{39}-8x_{39}^{2}-4x_{40}^{2}+x_{41}
\end{align*}
% END_RESIDUAL 26

% BEGIN_RESIDUAL 27
\begin{align*}
r_{27}={}&-16x_{8}^{2}+24x_{8}x_{10}-16x_{8}x_{19}+16x_{8}x_{27}+x_{10}-9x_{10}^{2}+12x_{10}x_{19}-12x_{10}x_{27}-4x_{19}^{2}+8x_{19}x_{27}\\
&-4x_{27}^{2}+x_{42}
\end{align*}
% END_RESIDUAL 27

% BEGIN_RESIDUAL 28
\begin{align*}
r_{28}={}&-1+6x_{9}-4x_{9}^{2}+8x_{9}x_{10}-8x_{9}x_{23}-6x_{10}-4x_{10}^{2}+8x_{10}x_{23}+6x_{23}-4x_{23}^{2}+x_{43}
\end{align*}
% END_RESIDUAL 28

% BEGIN_RESIDUAL 29
\begin{align*}
r_{29}={}&-1-x_{10}-x_{10}^{2}-4x_{10}x_{43}-4x_{43}^{2}+x_{44}
\end{align*}
% END_RESIDUAL 29

% BEGIN_RESIDUAL 30
\begin{align*}
r_{30}={}&1+x_{9}-x_{9}^{2}-4x_{9}x_{44}-4x_{44}^{2}+x_{45}
\end{align*}
% END_RESIDUAL 30

% BEGIN_RESIDUAL 31
\begin{align*}
r_{31}={}&-1-6x_{45}-4x_{45}^{2}+x_{46}
\end{align*}
% END_RESIDUAL 31

% BEGIN_RESIDUAL 32
\begin{align*}
r_{32}={}&-1-x_{10}-x_{10}^{2}-4x_{10}x_{46}-4x_{46}^{2}+x_{47}
\end{align*}
% END_RESIDUAL 32

% BEGIN_RESIDUAL 33
\begin{align*}
r_{33}={}&2+x_{9}-x_{9}^{2}-4x_{9}x_{47}-4x_{47}^{2}+x_{48}
\end{align*}
% END_RESIDUAL 33

% BEGIN_RESIDUAL 34
\begin{align*}
r_{34}={}&-2x_{42}-4x_{42}^{2}-8x_{42}x_{48}-4x_{48}^{2}+x_{49}
\end{align*}
% END_RESIDUAL 34

% BEGIN_RESIDUAL 35
\begin{align*}
r_{35}={}&-9+7x_{8}-x_{8}^{2}-4x_{8}x_{49}+12x_{49}-4x_{49}^{2}+x_{50}
\end{align*}
% END_RESIDUAL 35

% BEGIN_RESIDUAL 36
\begin{align*}
r_{36}={}&-1-6x_{50}-4x_{50}^{2}+x_{51}
\end{align*}
% END_RESIDUAL 36

% BEGIN_RESIDUAL 37
\begin{align*}
r_{37}={}&-x_{10}-x_{10}^{2}-4x_{10}x_{51}-4x_{51}^{2}+x_{52}
\end{align*}
% END_RESIDUAL 37

% BEGIN_RESIDUAL 38
\begin{align*}
r_{38}={}&-x_{8}-x_{8}^{2}-4x_{8}x_{52}-4x_{52}-4x_{52}^{2}+x_{53}
\end{align*}
% END_RESIDUAL 38

% BEGIN_RESIDUAL 39
\begin{align*}
r_{39}={}&-10-x_{9}+x_{10}-x_{23}-14x_{53}-4x_{53}^{2}+x_{54}
\end{align*}
% END_RESIDUAL 39

% BEGIN_RESIDUAL 40
\begin{align*}
r_{40}={}&1-2x_{9}+2x_{10}-2x_{23}+4x_{54}-4x_{54}^{2}+x_{55}
\end{align*}
% END_RESIDUAL 40

% BEGIN_RESIDUAL 41
\begin{align*}
r_{41}={}&-10+x_{9}-x_{9}^{2}-4x_{9}x_{55}-14x_{48}-4x_{48}^{2}-4x_{55}^{2}+x_{56}
\end{align*}
% END_RESIDUAL 41

% BEGIN_RESIDUAL 42
\begin{align*}
r_{42}={}&21+28x_{48}+8x_{48}^{2}-2x_{56}-4x_{56}^{2}+x_{57}
\end{align*}
% END_RESIDUAL 42

% BEGIN_RESIDUAL 43
\begin{align*}
r_{43}={}&15x_{9}-36x_{9}^{2}+12x_{9}x_{10}+24x_{9}x_{19}-x_{10}-x_{10}^{2}-4x_{10}x_{19}-5x_{19}-4x_{19}^{2}+x_{58}
\end{align*}
% END_RESIDUAL 43

% BEGIN_RESIDUAL 44
\begin{align*}
r_{44}={}&1+6x_{9}+x_{10}-x_{10}^{2}-4x_{10}x_{57}-2x_{19}-4x_{57}^{2}+4x_{58}-4x_{58}^{2}+x_{59}
\end{align*}
% END_RESIDUAL 44

% BEGIN_RESIDUAL 45
\begin{align*}
r_{45}={}&-2-2x_{10}+2x_{10}^{2}+8x_{10}x_{57}+8x_{57}^{2}-2x_{59}-4x_{59}^{2}+x_{60}
\end{align*}
% END_RESIDUAL 45

% BEGIN_RESIDUAL 46
\begin{align*}
r_{46}={}&-2-4x_{60}-4x_{60}^{2}+x_{61}
\end{align*}
% END_RESIDUAL 46

% BEGIN_RESIDUAL 47
\begin{align*}
r_{47}={}&-x_{3}x_{13}+x_{62}
\end{align*}
% END_RESIDUAL 47

% BEGIN_RESIDUAL 48
\begin{align*}
r_{48}={}&x_{12}-x_{13}-2x_{22}x_{62}-x_{63}
\end{align*}
% END_RESIDUAL 48

% BEGIN_RESIDUAL 49
\begin{align*}
r_{49}={}&12x_{10}x_{62}+x_{12}-x_{13}-2x_{26}x_{62}-x_{64}
\end{align*}
% END_RESIDUAL 49

% BEGIN_RESIDUAL 50
\begin{align*}
r_{50}={}&x_{12}-x_{13}-2x_{28}x_{62}-6x_{62}-x_{65}
\end{align*}
% END_RESIDUAL 50

% BEGIN_RESIDUAL 51
\begin{align*}
r_{51}={}&-12x_{10}x_{62}+x_{12}-x_{13}-2x_{30}x_{62}-x_{66}
\end{align*}
% END_RESIDUAL 51

% BEGIN_RESIDUAL 52
\begin{align*}
r_{52}={}&x_{12}-x_{13}-2x_{32}x_{62}+6x_{62}-x_{67}
\end{align*}
% END_RESIDUAL 52

% BEGIN_RESIDUAL 53
\begin{align*}
r_{53}={}&x_{12}-x_{13}-4x_{34}x_{62}-2x_{35}x_{62}-x_{68}
\end{align*}
% END_RESIDUAL 53

% BEGIN_RESIDUAL 54
\begin{align*}
r_{54}={}&x_{12}-x_{13}-2x_{38}x_{62}-x_{69}
\end{align*}
% END_RESIDUAL 54

% BEGIN_RESIDUAL 55
\begin{align*}
r_{55}={}&x_{12}-x_{13}-2x_{41}x_{62}-4x_{62}-x_{70}
\end{align*}
% END_RESIDUAL 55

% BEGIN_RESIDUAL 56
\begin{align*}
r_{56}={}&-x_{8}x_{9}+x_{8}x_{10}+x_{9}+x_{9}x_{10}-x_{10}-x_{10}^{2}+x_{71}
\end{align*}
% END_RESIDUAL 56

% BEGIN_RESIDUAL 57
\begin{align*}
r_{57}={}&1+2x_{11}+x_{11}^{2}-x_{71}^{2}+x_{72}
\end{align*}
% END_RESIDUAL 57

% BEGIN_RESIDUAL 58
\begin{align*}
r_{58}={}&-x_{12}^{2}+2x_{12}x_{13}+2x_{12}x_{74}-x_{13}^{2}-2x_{13}x_{74}-x_{72}^{2}+x_{73}-x_{74}^{2}
\end{align*}
% END_RESIDUAL 58

% BEGIN_RESIDUAL 59
\begin{align*}
r_{59}={}&-4x_{60}x_{62}-2x_{61}x_{62}+2x_{62}+x_{74}
\end{align*}
% END_RESIDUAL 59

% BEGIN_RESIDUAL 60
\begin{align*}
r_{60}={}&-x_{9}x_{62}-x_{13}+x_{75}x_{76}+x_{77}
\end{align*}
% END_RESIDUAL 60

% BEGIN_RESIDUAL 61
\begin{align*}
r_{61}={}&-1+2x_{3}+6x_{3}x_{9}-2x_{3}x_{19}+x_{75}
\end{align*}
% END_RESIDUAL 61

% BEGIN_RESIDUAL 62
\begin{align*}
r_{62}={}&-x_{8}x_{211}-x_{12}+x_{76}+x_{211}
\end{align*}
% END_RESIDUAL 62

% BEGIN_RESIDUAL 63
\begin{align*}
r_{63}={}&-4x_{8}x_{62}-2x_{9}x_{62}+x_{9}x_{211}+x_{12}-x_{13}-2x_{19}x_{62}+x_{78}
\end{align*}
% END_RESIDUAL 63

% BEGIN_RESIDUAL 64
\begin{align*}
r_{64}={}&-x_{12}x_{63}+x_{13}x_{63}+x_{79}
\end{align*}
% END_RESIDUAL 64

% BEGIN_RESIDUAL 65
\begin{align*}
r_{65}={}&-x_{64}x_{79}+x_{80}
\end{align*}
% END_RESIDUAL 65

% BEGIN_RESIDUAL 66
\begin{align*}
r_{66}={}&-x_{65}x_{80}+x_{81}
\end{align*}
% END_RESIDUAL 66

% BEGIN_RESIDUAL 67
\begin{align*}
r_{67}={}&-x_{66}x_{81}+x_{82}
\end{align*}
% END_RESIDUAL 67

% BEGIN_RESIDUAL 68
\begin{align*}
r_{68}={}&-x_{67}x_{82}+x_{83}
\end{align*}
% END_RESIDUAL 68

% BEGIN_RESIDUAL 69
\begin{align*}
r_{69}={}&-x_{68}x_{83}+x_{84}
\end{align*}
% END_RESIDUAL 69

% BEGIN_RESIDUAL 70
\begin{align*}
r_{70}={}&-x_{69}x_{84}+x_{85}
\end{align*}
% END_RESIDUAL 70

% BEGIN_RESIDUAL 71
\begin{align*}
r_{71}={}&-x_{70}x_{85}+x_{86}
\end{align*}
% END_RESIDUAL 71

% BEGIN_RESIDUAL 72
\begin{align*}
r_{72}={}&-x_{77}x_{86}+x_{87}
\end{align*}
% END_RESIDUAL 72

% BEGIN_RESIDUAL 73
\begin{align*}
r_{73}={}&-x_{78}x_{87}+x_{88}
\end{align*}
% END_RESIDUAL 73

% BEGIN_RESIDUAL 74
\begin{align*}
r_{74}={}&-x_{88}^{2}+x_{89}
\end{align*}
% END_RESIDUAL 74

% BEGIN_RESIDUAL 75
\begin{align*}
r_{75}={}&-x_{2}-x_{2}x_{16}+x_{16}+x_{90}
\end{align*}
% END_RESIDUAL 75

% BEGIN_RESIDUAL 76
\begin{align*}
r_{76}={}&-x_{14}x_{90}+x_{91}
\end{align*}
% END_RESIDUAL 76

% BEGIN_RESIDUAL 77
\begin{align*}
r_{77}={}&-2x_{1}^{2}+2x_{1}x_{12}+2x_{1}x_{13}+2x_{1}x_{91}+2x_{1}x_{92}-x_{12}^{2}-2x_{12}x_{91}-x_{13}^{2}-2x_{13}x_{92}+x_{17}x_{93}\\
&-x_{73}x_{89}-x_{91}^{2}-x_{92}^{2}
\end{align*}
% END_RESIDUAL 77

% BEGIN_RESIDUAL 78
\begin{align*}
r_{78}={}&-x_{2}x_{15}+x_{15}-x_{15}x_{90}+x_{92}
\end{align*}
% END_RESIDUAL 78

% BEGIN_RESIDUAL 79
\begin{align*}
r_{79}={}&-x_{96}-4x_{96}x_{209}+x_{98}
\end{align*}
% END_RESIDUAL 79

% BEGIN_RESIDUAL 80
\begin{align*}
r_{80}={}&-x_{4}x_{119}+x_{94}x_{201}-x_{97}-x_{98}x_{119}
\end{align*}
% END_RESIDUAL 80

% BEGIN_RESIDUAL 81
\begin{align*}
r_{81}={}&x_{99}x_{201}-x_{100}x_{122}-x_{101}
\end{align*}
% END_RESIDUAL 81

% BEGIN_RESIDUAL 82
\begin{align*}
r_{82}={}&x_{102}x_{201}-x_{103}x_{122}-x_{104}
\end{align*}
% END_RESIDUAL 82

% BEGIN_RESIDUAL 83
\begin{align*}
r_{83}={}&x_{105}x_{201}-x_{106}x_{122}-x_{107}
\end{align*}
% END_RESIDUAL 83

% BEGIN_RESIDUAL 84
\begin{align*}
r_{84}={}&x_{108}x_{201}-x_{109}x_{122}-x_{110}
\end{align*}
% END_RESIDUAL 84

% BEGIN_RESIDUAL 85
\begin{align*}
r_{85}={}&x_{111}x_{202}-x_{113}x_{122}-x_{114}
\end{align*}
% END_RESIDUAL 85

% BEGIN_RESIDUAL 86
\begin{align*}
r_{86}={}&x_{115}x_{202}-x_{116}x_{122}-x_{117}
\end{align*}
% END_RESIDUAL 86

% BEGIN_RESIDUAL 87
\begin{align*}
r_{87}={}&-x_{17}x_{119}+x_{122}
\end{align*}
% END_RESIDUAL 87

% BEGIN_RESIDUAL 88
\begin{align*}
r_{88}={}&x_{118}-x_{120}x_{122}-x_{121}
\end{align*}
% END_RESIDUAL 88

% BEGIN_RESIDUAL 89
\begin{align*}
r_{89}={}&-x_{122}x_{124}+x_{123}-x_{125}
\end{align*}
% END_RESIDUAL 89

% BEGIN_RESIDUAL 90
\begin{align*}
r_{90}={}&-x_{122}x_{127}+x_{126}x_{201}-x_{128}
\end{align*}
% END_RESIDUAL 90

% BEGIN_RESIDUAL 91
\begin{align*}
r_{91}={}&-x_{122}x_{130}+x_{129}x_{201}-x_{131}
\end{align*}
% END_RESIDUAL 91

% BEGIN_RESIDUAL 92
\begin{align*}
r_{92}={}&x_{132}-2x_{132}x_{134}+x_{135}
\end{align*}
% END_RESIDUAL 92

% BEGIN_RESIDUAL 93
\begin{align*}
r_{93}={}&-x_{132}x_{135}-x_{135}+x_{136}
\end{align*}
% END_RESIDUAL 93

% BEGIN_RESIDUAL 94
\begin{align*}
r_{94}={}&2x_{133}+x_{133}^{2}+x_{135}^{2}-x_{136}^{2}
\end{align*}
% END_RESIDUAL 94

% BEGIN_RESIDUAL 95
\begin{align*}
r_{95}={}&-x_{133}x_{138}-x_{138}+x_{139}
\end{align*}
% END_RESIDUAL 95

% BEGIN_RESIDUAL 96
\begin{align*}
r_{96}={}&-1+x_{137}^{2}+2x_{137}x_{139}+x_{138}^{2}
\end{align*}
% END_RESIDUAL 96

% BEGIN_RESIDUAL 97
\begin{align*}
r_{97}={}&1-2x_{0}+x_{0}^{2}-2x_{0}x_{133}+x_{140}
\end{align*}
% END_RESIDUAL 97

% BEGIN_RESIDUAL 98
\begin{align*}
r_{98}={}&-x_{0}x_{138}+x_{3}-x_{137}+x_{140}x_{141}
\end{align*}
% END_RESIDUAL 98

% BEGIN_RESIDUAL 99
\begin{align*}
r_{99}={}&1-2x_{4}+x_{4}^{2}-2x_{4}x_{133}+x_{142}
\end{align*}
% END_RESIDUAL 99

% BEGIN_RESIDUAL 100
\begin{align*}
r_{100}={}&-x_{4}x_{174}+x_{17}+x_{142}x_{143}-x_{173}
\end{align*}
% END_RESIDUAL 100

% BEGIN_RESIDUAL 101
\begin{align*}
r_{101}={}&-1+2x_{133}x_{185}-x_{144}+2x_{185}-x_{185}^{2}
\end{align*}
% END_RESIDUAL 101

% BEGIN_RESIDUAL 102
\begin{align*}
r_{102}={}&1+x_{144}x_{145}-x_{173}-x_{174}x_{185}+4x_{209}
\end{align*}
% END_RESIDUAL 102

% BEGIN_RESIDUAL 103
\begin{align*}
r_{103}={}&-x_{133}x_{147}-x_{147}+x_{148}
\end{align*}
% END_RESIDUAL 103

% BEGIN_RESIDUAL 104
\begin{align*}
r_{104}={}&-1+x_{146}^{2}+2x_{146}x_{148}+x_{147}^{2}
\end{align*}
% END_RESIDUAL 104

% BEGIN_RESIDUAL 105
\begin{align*}
r_{105}={}&1+x_{133}-x_{133}x_{150}-x_{150}+x_{151}
\end{align*}
% END_RESIDUAL 105

% BEGIN_RESIDUAL 106
\begin{align*}
r_{106}={}&x_{149}^{2}+2x_{149}x_{151}-2x_{150}+x_{150}^{2}
\end{align*}
% END_RESIDUAL 106

% BEGIN_RESIDUAL 107
\begin{align*}
r_{107}={}&-x_{133}x_{153}-x_{153}+x_{154}
\end{align*}
% END_RESIDUAL 107

% BEGIN_RESIDUAL 108
\begin{align*}
r_{108}={}&-1+x_{152}^{2}+2x_{152}x_{154}+x_{153}^{2}
\end{align*}
% END_RESIDUAL 108

% BEGIN_RESIDUAL 109
\begin{align*}
r_{109}={}&-x_{133}x_{156}-x_{156}+x_{157}
\end{align*}
% END_RESIDUAL 109

% BEGIN_RESIDUAL 110
\begin{align*}
r_{110}={}&-1+x_{155}^{2}+2x_{155}x_{157}+x_{156}^{2}
\end{align*}
% END_RESIDUAL 110

% BEGIN_RESIDUAL 111
\begin{align*}
r_{111}={}&-x_{133}x_{159}-x_{159}+x_{160}
\end{align*}
% END_RESIDUAL 111

% BEGIN_RESIDUAL 112
\begin{align*}
r_{112}={}&-1+x_{158}^{2}+2x_{158}x_{160}+x_{159}^{2}
\end{align*}
% END_RESIDUAL 112

% BEGIN_RESIDUAL 113
\begin{align*}
r_{113}={}&-x_{133}x_{162}-x_{162}+x_{163}
\end{align*}
% END_RESIDUAL 113

% BEGIN_RESIDUAL 114
\begin{align*}
r_{114}={}&-1+x_{161}^{2}+2x_{161}x_{163}+x_{162}^{2}
\end{align*}
% END_RESIDUAL 114

% BEGIN_RESIDUAL 115
\begin{align*}
r_{115}={}&-x_{133}x_{165}-x_{165}+x_{166}
\end{align*}
% END_RESIDUAL 115

% BEGIN_RESIDUAL 116
\begin{align*}
r_{116}={}&-1+x_{164}^{2}+2x_{164}x_{166}+x_{165}^{2}
\end{align*}
% END_RESIDUAL 116

% BEGIN_RESIDUAL 117
\begin{align*}
r_{117}={}&-x_{133}x_{168}-x_{168}+x_{169}
\end{align*}
% END_RESIDUAL 117

% BEGIN_RESIDUAL 118
\begin{align*}
r_{118}={}&-1+x_{167}^{2}+2x_{167}x_{169}+x_{168}^{2}
\end{align*}
% END_RESIDUAL 118

% BEGIN_RESIDUAL 119
\begin{align*}
r_{119}={}&-x_{133}x_{171}-x_{171}+x_{172}
\end{align*}
% END_RESIDUAL 119

% BEGIN_RESIDUAL 120
\begin{align*}
r_{120}={}&-1+x_{170}^{2}+2x_{170}x_{172}+x_{171}^{2}
\end{align*}
% END_RESIDUAL 120

% BEGIN_RESIDUAL 121
\begin{align*}
r_{121}={}&1-x_{133}x_{174}-x_{174}+x_{175}
\end{align*}
% END_RESIDUAL 121

% BEGIN_RESIDUAL 122
\begin{align*}
r_{122}={}&-1+2x_{173}+x_{173}^{2}+2x_{173}x_{175}+x_{174}^{2}
\end{align*}
% END_RESIDUAL 122

% BEGIN_RESIDUAL 123
\begin{align*}
r_{123}={}&-x_{133}x_{177}-x_{177}+x_{178}
\end{align*}
% END_RESIDUAL 123

% BEGIN_RESIDUAL 124
\begin{align*}
r_{124}={}&-1+x_{176}^{2}+2x_{176}x_{178}+x_{177}^{2}
\end{align*}
% END_RESIDUAL 124

% BEGIN_RESIDUAL 125
\begin{align*}
r_{125}={}&-x_{133}x_{180}-x_{180}+x_{181}
\end{align*}
% END_RESIDUAL 125

% BEGIN_RESIDUAL 126
\begin{align*}
r_{126}={}&-1+x_{179}^{2}+2x_{179}x_{181}+x_{180}^{2}
\end{align*}
% END_RESIDUAL 126

% BEGIN_RESIDUAL 127
\begin{align*}
r_{127}={}&-x_{133}x_{183}-x_{183}+x_{184}
\end{align*}
% END_RESIDUAL 127

% BEGIN_RESIDUAL 128
\begin{align*}
r_{128}={}&-1+x_{182}^{2}+2x_{182}x_{184}+x_{183}^{2}
\end{align*}
% END_RESIDUAL 128

% BEGIN_RESIDUAL 129
\begin{align*}
r_{129}={}&-8-2x_{0}^{2}-2x_{0}x_{3}-x_{1}^{2}-x_{2}^{2}-2x_{3}^{2}-x_{4}^{2}+2x_{8}-26x_{8}^{2}+30x_{8}x_{9}-40x_{8}x_{19}-20x_{8}x_{22}\\
&-20x_{8}x_{26}-20x_{8}x_{28}-20x_{8}x_{30}-20x_{8}x_{32}-40x_{8}x_{34}-20x_{8}x_{35}-20x_{8}x_{38}-20x_{8}x_{41}-40x_{8}x_{60}\\
&-20x_{8}x_{61}-10x_{9}^{2}+24x_{9}x_{19}+12x_{9}x_{22}+12x_{9}x_{26}+12x_{9}x_{28}+12x_{9}x_{30}+12x_{9}x_{32}+24x_{9}x_{34}\\
&+12x_{9}x_{35}+12x_{9}x_{38}+12x_{9}x_{41}+24x_{9}x_{60}+12x_{9}x_{61}-x_{10}^{2}-x_{11}^{2}-x_{12}^{2}-x_{13}^{2}-x_{14}^{2}\\
&-x_{15}^{2}-x_{16}^{2}-x_{17}^{2}-16x_{19}^{2}-16x_{19}x_{22}-16x_{19}x_{26}-16x_{19}x_{28}-16x_{19}x_{30}-16x_{19}x_{32}\\
&-32x_{19}x_{34}-16x_{19}x_{35}-16x_{19}x_{38}-16x_{19}x_{41}-32x_{19}x_{60}-16x_{19}x_{61}-4x_{22}^{2}-8x_{22}x_{26}-8x_{22}x_{28}\\
&-8x_{22}x_{30}-8x_{22}x_{32}-16x_{22}x_{34}-8x_{22}x_{35}-8x_{22}x_{38}-8x_{22}x_{41}-16x_{22}x_{60}-8x_{22}x_{61}-4x_{26}^{2}\\
&-8x_{26}x_{28}-8x_{26}x_{30}-8x_{26}x_{32}-16x_{26}x_{34}-8x_{26}x_{35}-8x_{26}x_{38}-8x_{26}x_{41}-16x_{26}x_{60}-8x_{26}x_{61}\\
&-4x_{28}^{2}-8x_{28}x_{30}-8x_{28}x_{32}-16x_{28}x_{34}-8x_{28}x_{35}-8x_{28}x_{38}-8x_{28}x_{41}-16x_{28}x_{60}-8x_{28}x_{61}\\
&-4x_{30}^{2}-8x_{30}x_{32}-16x_{30}x_{34}-8x_{30}x_{35}-8x_{30}x_{38}-8x_{30}x_{41}-16x_{30}x_{60}-8x_{30}x_{61}-4x_{32}^{2}\\
&-16x_{32}x_{34}-8x_{32}x_{35}-8x_{32}x_{38}-8x_{32}x_{41}-16x_{32}x_{60}-8x_{32}x_{61}-16x_{34}^{2}-16x_{34}x_{35}-16x_{34}x_{38}\\
&-16x_{34}x_{41}-32x_{34}x_{60}-16x_{34}x_{61}-4x_{35}^{2}-8x_{35}x_{38}-8x_{35}x_{41}-16x_{35}x_{60}-8x_{35}x_{61}-4x_{38}^{2}\\
&-8x_{38}x_{41}-16x_{38}x_{60}-8x_{38}x_{61}-4x_{41}^{2}-16x_{41}x_{60}-8x_{41}x_{61}-16x_{60}^{2}-16x_{60}x_{61}-4x_{61}^{2}\\
&-x_{94}^{2}-x_{95}^{2}-x_{99}^{2}-x_{102}^{2}-x_{105}^{2}-x_{108}^{2}-x_{111}^{2}-x_{112}^{2}-x_{115}^{2}-x_{118}^{2}-x_{119}^{2}\\
&-x_{123}^{2}-x_{126}^{2}-x_{129}^{2}+x_{132}-x_{185}^{2}-24x_{209}-32x_{209}^{2}
\end{align*}
% END_RESIDUAL 129

% BEGIN_RESIDUAL 130
\begin{align*}
r_{130}={}&-x_{4}-x_{4}x_{8}-x_{4}x_{9}-x_{4}x_{10}-x_{4}x_{11}-x_{4}x_{12}-x_{4}x_{13}-x_{4}x_{14}-x_{4}x_{15}-x_{4}x_{16}+x_{185}
\end{align*}
% END_RESIDUAL 130

% BEGIN_RESIDUAL 131
\begin{align*}
r_{131}={}&-16+x_{0}-x_{186}^{2}
\end{align*}
% END_RESIDUAL 131

% BEGIN_RESIDUAL 132
\begin{align*}
r_{132}={}&-x_{133}x_{188}-x_{188}+x_{197}
\end{align*}
% END_RESIDUAL 132

% BEGIN_RESIDUAL 133
\begin{align*}
r_{133}={}&-1+x_{187}^{2}+x_{188}^{2}-x_{197}^{2}
\end{align*}
% END_RESIDUAL 133

% BEGIN_RESIDUAL 134
\begin{align*}
r_{134}={}&-x_{190}-x_{190}x_{191}+x_{198}
\end{align*}
% END_RESIDUAL 134

% BEGIN_RESIDUAL 135
\begin{align*}
r_{135}={}&-1+x_{189}^{2}+x_{190}^{2}-x_{198}^{2}
\end{align*}
% END_RESIDUAL 135

% BEGIN_RESIDUAL 136
\begin{align*}
r_{136}={}&-x_{183}^{2}+x_{199}
\end{align*}
% END_RESIDUAL 136

% BEGIN_RESIDUAL 137
\begin{align*}
r_{137}={}&x_{188}-2x_{193}x_{199}-x_{199}
\end{align*}
% END_RESIDUAL 137

% BEGIN_RESIDUAL 138
\begin{align*}
r_{138}={}&-4x_{188}-8x_{188}x_{192}+x_{191}
\end{align*}
% END_RESIDUAL 138

% BEGIN_RESIDUAL 139
\begin{align*}
r_{139}={}&-x_{133}-x_{187}x_{194}+x_{191}
\end{align*}
% END_RESIDUAL 139

% BEGIN_RESIDUAL 140
\begin{align*}
r_{140}={}&-x_{182}-x_{184}-x_{187}x_{195}+x_{189}
\end{align*}
% END_RESIDUAL 140

% BEGIN_RESIDUAL 141
\begin{align*}
r_{141}={}&-5x_{8}+3x_{9}-4x_{19}-2x_{22}-2x_{26}-2x_{28}-2x_{30}-2x_{32}-4x_{34}-2x_{35}-2x_{38}-2x_{41}-4x_{60}-2x_{61}-4x_{183}x_{196}\\
&+x_{190}
\end{align*}
% END_RESIDUAL 141

% BEGIN_RESIDUAL 142
\begin{align*}
r_{142}={}&26+2x_{0}^{2}+2x_{0}x_{3}-2x_{0}x_{138}-2x_{0}x_{147}+x_{1}^{2}-2x_{1}x_{174}+x_{2}^{2}-2x_{2}x_{165}+x_{3}^{2}-2x_{3}x_{147}\\
&+26x_{8}^{2}-30x_{8}x_{9}+40x_{8}x_{19}+20x_{8}x_{22}+20x_{8}x_{26}+20x_{8}x_{28}+20x_{8}x_{30}+20x_{8}x_{32}+40x_{8}x_{34}\\
&+20x_{8}x_{35}+20x_{8}x_{38}+20x_{8}x_{41}+40x_{8}x_{60}+20x_{8}x_{61}-2x_{8}x_{150}-10x_{8}x_{183}+10x_{9}^{2}-24x_{9}x_{19}\\
&-12x_{9}x_{22}-12x_{9}x_{26}-12x_{9}x_{28}-12x_{9}x_{30}-12x_{9}x_{32}-24x_{9}x_{34}-12x_{9}x_{35}-12x_{9}x_{38}-12x_{9}x_{41}\\
&-24x_{9}x_{60}-12x_{9}x_{61}-2x_{9}x_{153}+6x_{9}x_{183}+x_{10}^{2}-2x_{10}x_{156}+x_{11}^{2}-2x_{11}x_{159}+x_{12}^{2}\\
&-2x_{12}x_{162}+x_{13}^{2}-2x_{13}x_{168}+x_{14}^{2}-2x_{14}x_{171}+x_{15}^{2}-2x_{15}x_{177}+x_{16}^{2}-2x_{16}x_{180}\\
&+16x_{19}^{2}+16x_{19}x_{22}+16x_{19}x_{26}+16x_{19}x_{28}+16x_{19}x_{30}+16x_{19}x_{32}+32x_{19}x_{34}+16x_{19}x_{35}\\
&+16x_{19}x_{38}+16x_{19}x_{41}+32x_{19}x_{60}+16x_{19}x_{61}-8x_{19}x_{183}+4x_{22}^{2}+8x_{22}x_{26}+8x_{22}x_{28}+8x_{22}x_{30}\\
&+8x_{22}x_{32}+16x_{22}x_{34}+8x_{22}x_{35}+8x_{22}x_{38}+8x_{22}x_{41}+16x_{22}x_{60}+8x_{22}x_{61}-4x_{22}x_{183}+4x_{26}^{2}\\
&+8x_{26}x_{28}+8x_{26}x_{30}+8x_{26}x_{32}+16x_{26}x_{34}+8x_{26}x_{35}+8x_{26}x_{38}+8x_{26}x_{41}+16x_{26}x_{60}+8x_{26}x_{61}\\
&-4x_{26}x_{183}+4x_{28}^{2}+8x_{28}x_{30}+8x_{28}x_{32}+16x_{28}x_{34}+8x_{28}x_{35}+8x_{28}x_{38}+8x_{28}x_{41}+16x_{28}x_{60}\\
&+8x_{28}x_{61}-4x_{28}x_{183}+4x_{30}^{2}+8x_{30}x_{32}+16x_{30}x_{34}+8x_{30}x_{35}+8x_{30}x_{38}+8x_{30}x_{41}+16x_{30}x_{60}\\
&+8x_{30}x_{61}-4x_{30}x_{183}+4x_{32}^{2}+16x_{32}x_{34}+8x_{32}x_{35}+8x_{32}x_{38}+8x_{32}x_{41}+16x_{32}x_{60}+8x_{32}x_{61}\\
&-4x_{32}x_{183}+16x_{34}^{2}+16x_{34}x_{35}+16x_{34}x_{38}+16x_{34}x_{41}+32x_{34}x_{60}+16x_{34}x_{61}-8x_{34}x_{183}+4x_{35}^{2}\\
&+8x_{35}x_{38}+8x_{35}x_{41}+16x_{35}x_{60}+8x_{35}x_{61}-4x_{35}x_{183}+4x_{38}^{2}+8x_{38}x_{41}+16x_{38}x_{60}+8x_{38}x_{61}\\
&-4x_{38}x_{183}+4x_{41}^{2}+16x_{41}x_{60}+8x_{41}x_{61}-4x_{41}x_{183}+16x_{60}^{2}+16x_{60}x_{61}-8x_{60}x_{183}+4x_{61}^{2}\\
&-4x_{61}x_{183}-x_{133}x_{200}-2x_{137}-2x_{139}+x_{139}^{2}-2x_{146}-2x_{148}+x_{148}^{2}-2x_{149}+2x_{150}-2x_{151}+x_{151}^{2}\\
&-2x_{152}-2x_{154}+x_{154}^{2}-2x_{155}-2x_{157}+x_{157}^{2}-2x_{158}-2x_{160}+x_{160}^{2}-2x_{161}-2x_{163}+x_{163}^{2}-2x_{164}\\
&-2x_{166}+x_{166}^{2}-2x_{167}-2x_{169}+x_{169}^{2}-2x_{170}-2x_{172}+x_{172}^{2}-2x_{173}+x_{175}^{2}-2x_{176}-2x_{178}\\
&+x_{178}^{2}-2x_{179}-2x_{181}+x_{181}^{2}-2x_{182}-2x_{184}+x_{184}^{2}
\end{align*}
% END_RESIDUAL 142

% BEGIN_RESIDUAL 143
\begin{align*}
r_{143}={}&-x_{95}x_{201}+x_{119}
\end{align*}
% END_RESIDUAL 143

% BEGIN_RESIDUAL 144
\begin{align*}
r_{144}={}&-x_{112}x_{202}+x_{119}
\end{align*}
% END_RESIDUAL 144

% BEGIN_RESIDUAL 145
\begin{align*}
r_{145}={}&-1+x_{2}-x_{3}-2x_{3}x_{210}
\end{align*}
% END_RESIDUAL 145

% BEGIN_RESIDUAL 146
\begin{align*}
r_{146}={}&-2-25x_{8}^{2}+30x_{8}x_{9}-40x_{8}x_{19}-20x_{8}x_{22}-20x_{8}x_{26}-20x_{8}x_{28}-20x_{8}x_{30}-20x_{8}x_{32}-40x_{8}x_{34}\\
&-20x_{8}x_{35}-20x_{8}x_{38}-20x_{8}x_{41}-40x_{8}x_{60}-20x_{8}x_{61}-9x_{9}^{2}+24x_{9}x_{19}+12x_{9}x_{22}+12x_{9}x_{26}\\
&+12x_{9}x_{28}+12x_{9}x_{30}+12x_{9}x_{32}+24x_{9}x_{34}+12x_{9}x_{35}+12x_{9}x_{38}+12x_{9}x_{41}+24x_{9}x_{60}+12x_{9}x_{61}\\
&-16x_{19}^{2}-16x_{19}x_{22}-16x_{19}x_{26}-16x_{19}x_{28}-16x_{19}x_{30}-16x_{19}x_{32}-32x_{19}x_{34}-16x_{19}x_{35}\\
&-16x_{19}x_{38}-16x_{19}x_{41}-32x_{19}x_{60}-16x_{19}x_{61}-4x_{22}^{2}-8x_{22}x_{26}-8x_{22}x_{28}-8x_{22}x_{30}-8x_{22}x_{32}\\
&-16x_{22}x_{34}-8x_{22}x_{35}-8x_{22}x_{38}-8x_{22}x_{41}-16x_{22}x_{60}-8x_{22}x_{61}-4x_{26}^{2}-8x_{26}x_{28}-8x_{26}x_{30}\\
&-8x_{26}x_{32}-16x_{26}x_{34}-8x_{26}x_{35}-8x_{26}x_{38}-8x_{26}x_{41}-16x_{26}x_{60}-8x_{26}x_{61}-4x_{28}^{2}-8x_{28}x_{30}\\
&-8x_{28}x_{32}-16x_{28}x_{34}-8x_{28}x_{35}-8x_{28}x_{38}-8x_{28}x_{41}-16x_{28}x_{60}-8x_{28}x_{61}-4x_{30}^{2}-8x_{30}x_{32}\\
&-16x_{30}x_{34}-8x_{30}x_{35}-8x_{30}x_{38}-8x_{30}x_{41}-16x_{30}x_{60}-8x_{30}x_{61}-4x_{32}^{2}-16x_{32}x_{34}-8x_{32}x_{35}\\
&-8x_{32}x_{38}-8x_{32}x_{41}-16x_{32}x_{60}-8x_{32}x_{61}-16x_{34}^{2}-16x_{34}x_{35}-16x_{34}x_{38}-16x_{34}x_{41}-32x_{34}x_{60}\\
&-16x_{34}x_{61}-4x_{35}^{2}-8x_{35}x_{38}-8x_{35}x_{41}-16x_{35}x_{60}-8x_{35}x_{61}-4x_{38}^{2}-8x_{38}x_{41}-16x_{38}x_{60}\\
&-8x_{38}x_{61}-4x_{41}^{2}-16x_{41}x_{60}-8x_{41}x_{61}-16x_{60}^{2}-16x_{60}x_{61}-4x_{61}^{2}-x_{138}^{2}-x_{147}^{2}+2x_{150}\\
&-x_{150}^{2}-x_{153}^{2}-x_{156}^{2}-x_{159}^{2}-x_{162}^{2}-x_{165}^{2}-x_{168}^{2}-x_{171}^{2}-x_{174}^{2}-x_{177}^{2}\\
&-x_{180}^{2}+x_{183}^{2}-x_{203}^{2}-x_{204}^{2}-x_{205}-x_{205}^{2}
\end{align*}
% END_RESIDUAL 146

% BEGIN_RESIDUAL 147
\begin{align*}
r_{147}={}&-1-2x_{97}^{2}-2x_{101}^{2}-2x_{104}^{2}-2x_{107}^{2}-2x_{110}^{2}-2x_{114}^{2}-2x_{117}^{2}+x_{119}^{2}-2x_{121}^{2}-2x_{125}^{2}\\
&-2x_{128}^{2}-2x_{131}^{2}-x_{206}^{2}-x_{207}^{2}-x_{208}-x_{208}^{2}
\end{align*}
% END_RESIDUAL 147

% BEGIN_RESIDUAL 148
\begin{align*}
r_{148}={}&-x_{3}x_{12}+x_{211}
\end{align*}
% END_RESIDUAL 148

% BEGIN_RESIDUAL 149
\begin{align*}
r_{149}={}&-2x_{94}x_{147}+2x_{99}-2x_{99}x_{150}-2x_{102}x_{153}-2x_{105}x_{156}-2x_{108}x_{159}-2x_{111}x_{162}-2x_{115}x_{168}\\
&-2x_{118}x_{171}-2x_{123}x_{177}-2x_{126}x_{180}-2x_{129}x_{183}-2x_{133}x_{213}+2x_{146}x_{147}+2x_{147}x_{148}-2x_{149}\\
&+2x_{149}x_{150}+2x_{150}x_{151}-2x_{151}+2x_{152}x_{153}+2x_{153}x_{154}+2x_{155}x_{156}+2x_{156}x_{157}+2x_{158}x_{159}\\
&+2x_{159}x_{160}+2x_{161}x_{162}+2x_{162}x_{163}+2x_{167}x_{168}+2x_{168}x_{169}+2x_{170}x_{171}+2x_{171}x_{172}+2x_{176}x_{177}\\
&+2x_{177}x_{178}+2x_{179}x_{180}+2x_{180}x_{181}+2x_{182}x_{183}+2x_{183}x_{184}+4x_{209}x_{213}+x_{212}
\end{align*}
% END_RESIDUAL 149

% BEGIN_RESIDUAL 150
\begin{align*}
r_{150}={}&-1+x_{94}^{2}-2x_{94}x_{146}+x_{99}^{2}-2x_{99}x_{149}+x_{102}^{2}-2x_{102}x_{152}+x_{105}^{2}-2x_{105}x_{155}+x_{108}^{2}\\
&-2x_{108}x_{158}+x_{111}^{2}-2x_{111}x_{161}+x_{115}^{2}-2x_{115}x_{167}+x_{118}^{2}-2x_{118}x_{170}+x_{123}^{2}-2x_{123}x_{176}\\
&+x_{126}^{2}-2x_{126}x_{179}+x_{129}^{2}-2x_{129}x_{182}+x_{146}^{2}-x_{147}^{2}+x_{149}^{2}+2x_{150}-x_{150}^{2}+x_{152}^{2}\\
&-x_{153}^{2}+x_{155}^{2}-x_{156}^{2}+x_{158}^{2}-x_{159}^{2}+x_{161}^{2}-x_{162}^{2}+x_{167}^{2}-x_{168}^{2}+x_{170}^{2}\\
&-x_{171}^{2}+x_{176}^{2}-x_{177}^{2}+x_{179}^{2}-x_{180}^{2}+x_{182}^{2}-x_{183}^{2}-4x_{209}x_{212}-2x_{212}+x_{213}
\end{align*}
% END_RESIDUAL 150

% BEGIN_RESIDUAL 151
\begin{align*}
r_{151}={}&-2x_{95}x_{138}-2x_{112}x_{165}-2x_{119}x_{174}-2x_{133}x_{215}+2x_{137}x_{138}+2x_{138}x_{139}+2x_{164}x_{165}+2x_{165}x_{166}\\
&+2x_{173}x_{174}+2x_{174}+2x_{174}x_{175}+4x_{209}x_{215}+x_{214}-x_{215}
\end{align*}
% END_RESIDUAL 151

% BEGIN_RESIDUAL 152
\begin{align*}
r_{152}={}&x_{95}^{2}-2x_{95}x_{137}+x_{112}^{2}-2x_{112}x_{164}+x_{119}^{2}-2x_{119}x_{173}+x_{137}^{2}-x_{138}^{2}+x_{164}^{2}-x_{165}^{2}\\
&+x_{173}^{2}-x_{174}^{2}-4x_{209}x_{214}-x_{214}+x_{215}
\end{align*}
% END_RESIDUAL 152

\endgroup


\section{Complete monic-elimination tables}\label{app:pivots}
For each profile the listed pair $(r,v)$ means solve original row $r_r$
for original coordinate $x_v$. All coefficients used for solving are
$\pm1$. Use recursive acyclic substitution, not a change in quantifier
domain. Coordinates outside \eqref{eq:deletedvars} and outside the
listed pivots remain free; append one new integer per block part.
These are fixed lists, with no data-dependent iteration in the
definition of the final polynomial.
\subsection{143 integers, degree 8}
Block partition: $(0;\ 1;\ 2;\ 3;\ 4)$.
There are 68 deleted monic coordinates and 80 final residuals.
\begin{center}\small\begin{tabular}{rrrrrr}\toprule
$(71,86)$ & $(50,65)$ & $(79,98)$ & $(37,52)$ & $(64,79)$ & $(73,88)$\\
$(57,72)$ & $(53,68)$ & $(81,101)$ & $(83,107)$ & $(61,75)$ & $(52,67)$\\
$(84,110)$ & $(89,125)$ & $(54,69)$ & $(88,121)$ & $(18,33)$ & $(85,114)$\\
$(78,92)$ & $(41,56)$ & $(136,199)$ & $(30,45)$ & $(91,131)$ & $(90,128)$\\
$(51,66)$ & $(82,104)$ & $(86,117)$ & $(59,74)$ & $(55,70)$ & $(43,58)$\\
$(49,64)$ & $(35,50)$ & $(93,136)$ & $(39,54)$ & $(62,76)$ & $(14,29)$\\
$(76,91)$ & $(16,31)$ & $(32,47)$ & $(140,189)$ & $(134,198)$ & $(25,40)$\\
$(0,7)$ & $(132,197)$ & $(10,25)$ & $(22,37)$ & $(27,42)$ & $(28,43)$\\
$(6,21)$ & $(115,166)$ & $(111,160)$ & $(123,178)$ & $(107,154)$ & $(119,172)$\\
$(125,181)$ & $(95,139)$ & $(105,151)$ & $(103,148)$ & $(121,175)$ & $(109,157)$\\
$(117,169)$ & $(113,163)$ & $(127,184)$ & $(1,1)$ & $(3,16)$ & $(20,35)$\\
$(13,28)$ & $(143,119)$ &  &  &  & \\
\bottomrule\end{tabular}\end{center}
\subsection{119 integers, degree 16}
Block partition: $(0,1,2;\ 3,4)$.
There are 89 deleted monic coordinates and 56 final residuals.
\begin{center}\small\begin{tabular}{rrrrrr}\toprule
$(53,68)$ & $(138,191)$ & $(52,67)$ & $(64,79)$ & $(54,69)$ & $(0,7)$\\
$(136,199)$ & $(50,65)$ & $(62,76)$ & $(72,87)$ & $(49,64)$ & $(63,78)$\\
$(55,70)$ & $(67,82)$ & $(85,114)$ & $(90,128)$ & $(81,101)$ & $(82,104)$\\
$(89,125)$ & $(80,97)$ & $(86,117)$ & $(88,121)$ & $(84,110)$ & $(83,107)$\\
$(91,131)$ & $(10,21)$ & $(48,63)$ & $(29,44)$ & $(76,91)$ & $(78,92)$\\
$(74,89)$ & $(16,31)$ & $(132,197)$ & $(37,52)$ & $(32,47)$ & $(44,59)$\\
$(27,42)$ & $(93,136)$ & $(61,75)$ & $(117,169)$ & $(105,151)$ & $(119,172)$\\
$(109,157)$ & $(127,184)$ & $(107,154)$ & $(111,160)$ & $(125,181)$ & $(1,1)$\\
$(113,163)$ & $(103,148)$ & $(115,166)$ & $(123,178)$ & $(41,56)$ & $(18,33)$\\
$(14,29)$ & $(51,66)$ & $(46,61)$ & $(20,35)$ & $(23,38)$ & $(13,28)$\\
$(56,71)$ & $(3,16)$ & $(25,40)$ & $(121,175)$ & $(95,139)$ & $(39,54)$\\
$(71,86)$ & $(141,190)$ & $(87,122)$ & $(79,98)$ & $(92,135)$ & $(69,84)$\\
$(143,119)$ & $(65,80)$ & $(59,74)$ & $(134,198)$ & $(75,90)$ & $(28,43)$\\
$(43,58)$ & $(34,49)$ & $(31,46)$ & $(24,39)$ & $(36,51)$ & $(21,36)$\\
$(5,20)$ & $(58,73)$ & $(140,189)$ & $(148,211)$ & $(9,24)$ & \\
\bottomrule\end{tabular}\end{center}
\subsection{112 integers, degree 24}
Block partition: $(0,1,2,3,4)$.
There are 95 deleted monic coordinates and 49 final residuals.
\begin{center}\small\begin{tabular}{rrrrrr}\toprule
$(53,68)$ & $(138,191)$ & $(52,67)$ & $(64,79)$ & $(54,69)$ & $(0,7)$\\
$(136,199)$ & $(50,65)$ & $(62,76)$ & $(72,87)$ & $(49,64)$ & $(63,78)$\\
$(55,70)$ & $(67,82)$ & $(85,114)$ & $(90,128)$ & $(81,101)$ & $(82,104)$\\
$(89,125)$ & $(80,97)$ & $(86,117)$ & $(88,121)$ & $(84,110)$ & $(83,107)$\\
$(91,131)$ & $(10,21)$ & $(48,63)$ & $(29,44)$ & $(76,91)$ & $(78,92)$\\
$(74,89)$ & $(16,31)$ & $(132,197)$ & $(37,52)$ & $(32,47)$ & $(44,59)$\\
$(27,42)$ & $(93,136)$ & $(61,75)$ & $(41,56)$ & $(18,33)$ & $(14,29)$\\
$(51,66)$ & $(56,71)$ & $(69,84)$ & $(25,40)$ & $(59,74)$ & $(39,54)$\\
$(71,86)$ & $(141,190)$ & $(1,1)$ & $(3,16)$ & $(20,35)$ & $(13,28)$\\
$(23,38)$ & $(46,61)$ & $(87,122)$ & $(79,98)$ & $(92,135)$ & $(117,169)$\\
$(105,151)$ & $(121,175)$ & $(119,172)$ & $(109,157)$ & $(127,184)$ & $(107,154)$\\
$(111,160)$ & $(95,139)$ & $(125,181)$ & $(113,163)$ & $(103,148)$ & $(115,166)$\\
$(123,178)$ & $(143,119)$ & $(65,80)$ & $(134,198)$ & $(75,90)$ & $(28,43)$\\
$(43,58)$ & $(34,49)$ & $(31,46)$ & $(24,39)$ & $(36,51)$ & $(21,36)$\\
$(5,20)$ & $(58,73)$ & $(140,189)$ & $(11,26)$ & $(57,72)$ & $(148,211)$\\
$(145,2)$ & $(68,83)$ & $(60,77)$ & $(9,24)$ & $(137,188)$ & \\
\bottomrule\end{tabular}\end{center}
\subsection{107 integers, degree 32}
Block partition: $(0,1,2,3,4)$.
There are 100 deleted monic coordinates and 44 final residuals.
\begin{center}\small\begin{tabular}{rrrrrr}\toprule
$(53,68)$ & $(138,191)$ & $(52,67)$ & $(64,79)$ & $(54,69)$ & $(0,7)$\\
$(136,199)$ & $(50,65)$ & $(62,76)$ & $(72,87)$ & $(49,64)$ & $(63,78)$\\
$(55,70)$ & $(67,82)$ & $(85,114)$ & $(90,128)$ & $(81,101)$ & $(82,104)$\\
$(89,125)$ & $(80,97)$ & $(86,117)$ & $(88,121)$ & $(84,110)$ & $(83,107)$\\
$(91,131)$ & $(10,21)$ & $(48,63)$ & $(29,44)$ & $(76,91)$ & $(78,92)$\\
$(74,89)$ & $(16,31)$ & $(132,197)$ & $(37,52)$ & $(32,47)$ & $(44,59)$\\
$(27,42)$ & $(93,136)$ & $(61,75)$ & $(41,56)$ & $(18,33)$ & $(14,29)$\\
$(51,66)$ & $(56,71)$ & $(69,84)$ & $(25,40)$ & $(59,74)$ & $(39,54)$\\
$(71,86)$ & $(141,190)$ & $(1,1)$ & $(3,16)$ & $(20,35)$ & $(13,28)$\\
$(23,38)$ & $(46,61)$ & $(87,122)$ & $(79,98)$ & $(92,135)$ & $(117,169)$\\
$(105,151)$ & $(121,175)$ & $(119,172)$ & $(109,157)$ & $(127,184)$ & $(107,154)$\\
$(111,160)$ & $(95,139)$ & $(125,181)$ & $(113,163)$ & $(103,148)$ & $(115,166)$\\
$(123,178)$ & $(143,119)$ & $(65,80)$ & $(134,198)$ & $(75,90)$ & $(28,43)$\\
$(43,58)$ & $(34,49)$ & $(31,46)$ & $(24,39)$ & $(36,51)$ & $(131,0)$\\
$(21,36)$ & $(5,20)$ & $(42,57)$ & $(58,73)$ & $(140,189)$ & $(11,26)$\\
$(57,72)$ & $(148,211)$ & $(145,2)$ & $(68,83)$ & $(60,77)$ & $(9,24)$\\
$(137,188)$ & $(47,62)$ & $(8,23)$ & $(4,19)$ &  & \\
\bottomrule\end{tabular}\end{center}
\clearpage
\subsection{102 integers, degree 64}
Block partition: $(0,1,2,3,4)$.
There are 105 deleted monic coordinates and 39 final residuals.
\begin{center}\small\begin{tabular}{rrrrrr}\toprule
$(72,87)$ & $(70,85)$ & $(74,89)$ & $(68,83)$ & $(136,199)$ & $(65,80)$\\
$(79,98)$ & $(138,191)$ & $(62,76)$ & $(48,63)$ & $(0,7)$ & $(51,66)$\\
$(76,91)$ & $(78,92)$ & $(58,73)$ & $(54,69)$ & $(50,65)$ & $(53,68)$\\
$(61,75)$ & $(63,78)$ & $(55,70)$ & $(52,67)$ & $(49,64)$ & $(129,132)$\\
$(132,197)$ & $(84,110)$ & $(90,128)$ & $(81,101)$ & $(82,104)$ & $(89,125)$\\
$(91,131)$ & $(85,114)$ & $(83,107)$ & $(88,121)$ & $(86,117)$ & $(36,51)$\\
$(140,189)$ & $(141,190)$ & $(31,46)$ & $(34,49)$ & $(38,53)$ & $(29,44)$\\
$(40,55)$ & $(56,71)$ & $(6,21)$ & $(22,37)$ & $(44,59)$ & $(25,40)$\\
$(16,31)$ & $(14,29)$ & $(10,25)$ & $(18,33)$ & $(80,97)$ & $(87,122)$\\
$(148,211)$ & $(67,82)$ & $(64,79)$ & $(59,74)$ & $(75,90)$ & $(93,136)$\\
$(71,86)$ & $(134,198)$ & $(115,166)$ & $(103,148)$ & $(117,169)$ & $(113,163)$\\
$(123,178)$ & $(95,139)$ & $(107,154)$ & $(119,172)$ & $(125,181)$ & $(109,157)$\\
$(111,160)$ & $(127,184)$ & $(121,175)$ & $(105,151)$ & $(24,39)$ & $(42,57)$\\
$(43,58)$ & $(28,43)$ & $(27,42)$ & $(32,47)$ & $(57,72)$ & $(92,135)$\\
$(60,77)$ & $(66,81)$ & $(21,36)$ & $(9,24)$ & $(5,20)$ & $(137,188)$\\
$(8,23)$ & $(1,1)$ & $(3,16)$ & $(46,61)$ & $(20,35)$ & $(143,119)$\\
$(13,28)$ & $(131,0)$ & $(37,52)$ & $(47,62)$ & $(145,2)$ & $(130,185)$\\
$(41,56)$ & $(139,133)$ & $(4,19)$ &  &  & \\
\bottomrule\end{tabular}\end{center}


\begin{thebibliography}{9}
\bibitem{source172} ProveIt research package 172,
\emph{Pooled powers and quadratic projections for ZFC consistency},
25 October 2026. The signed arithmetic component, its ordinary proof,
source calibration and affine216 chart are retained in this release.
\url{https://github.com/VladimirReshetnikov/ProveIt/tree/ef0c534c1ae1e312c61a2a9a5a3f18f5516fd42b/Computability/DiophantineConsistency/con-zfc-diophantine-forms/172-pooled-powers-round21}
\bibitem{pool168} ProveIt research package 168,
\emph{Reciprocal pooling of distinct-base Pell certificates},
25 October 2026, especially \texttt{proofs/RECIPROCAL\_POOLING.md}.
\url{https://github.com/VladimirReshetnikov/ProveIt/tree/ef0c534c1ae1e312c61a2a9a5a3f18f5516fd42b/Computability/DiophantineConsistency/con-zfc-diophantine-forms/168-reciprocal-pooling}
\bibitem{pool158} ProveIt research package 158,
\emph{Pooling exact power certificates for Diophantine consistency},
25 October 2026.
\url{https://github.com/VladimirReshetnikov/ProveIt/tree/ef0c534c1ae1e312c61a2a9a5a3f18f5516fd42b/Computability/DiophantineConsistency/con-zfc-diophantine-forms/158-pooled-powers}
\bibitem{eliminate152} ProveIt research package 152,
\emph{Lucas corrections and elimination of arithmetic witnesses},
25 October 2026, higher-degree component.
\url{https://github.com/VladimirReshetnikov/ProveIt/tree/ef0c534c1ae1e312c61a2a9a5a3f18f5516fd42b/Computability/DiophantineConsistency/con-zfc-diophantine-forms/152-lucas-corrections}
\bibitem{corpus} Vladimir Reshetnikov, \emph{ProveIt: Con(ZFC) in
Diophantine form}, repository corpus and comparison README, inspected
at the pinned snapshot stated in Section~1.
\url{https://github.com/VladimirReshetnikov/ProveIt/tree/ef0c534c1ae1e312c61a2a9a5a3f18f5516fd42b/Computability/DiophantineConsistency/con-zfc-diophantine-forms}
\end{thebibliography}
\end{document}
